% ============================================================================
% SYNCED COPY: manuscript/main.tex is identical to manuscript.tex (2026-08-19, final).
% Edit manuscript.tex; do not edit this file directly. See .paper/review/top_journal_review_20260818.md
% ============================================================================
\documentclass[11pt,a4paper]{article}
\usepackage[margin=1in]{geometry}
\usepackage{newtxtext,newtxmath}
\usepackage{setspace}
\singlespacing
\usepackage{amsmath}
\usepackage{graphicx}
\usepackage[svgnames]{xcolor}
\usepackage{booktabs}
\usepackage{array}
\usepackage[hidelinks]{hyperref}
\usepackage{bookmark}
\usepackage{url}
\usepackage{tabularx}
\usepackage{ragged2e}
\usepackage{tcolorbox}
\usepackage{caption}
\raggedbottom
\usepackage{fancyhdr}
\usepackage{enumitem}
\usepackage{placeins}
\setlist{nosep}
\setlist[itemize]{topsep=0em,itemsep=0em,parsep=0em,leftmargin=1.5em}
\setlist[enumerate]{topsep=0em,itemsep=0em,parsep=0em,leftmargin=1.5em}

% --- Cell Reports Methods: dedicated Summary environment equivalent to abstract (FFC requires 150 words) ---
\newenvironment{summary}{\vspace{0.6em}\noindent\textbf{Summary}\par\vspace{0.3em}\noindent\ignorespaces}{\par\vspace{0.5em}}

\pagestyle{fancy}
\fancyhf{}
\fancyhead[L]{\small Self-Supervised Pretraining on Multi-Cancer Molecular Networks}
\fancyhead[R]{\small}
\fancyfoot[C]{\thepage}
\renewcommand{\headrulewidth}{0.4pt}

\title{\textbf{Self-Supervised Pretraining on Multi-Cancer Molecular Networks:\\[-0.2em]Task Design Principles and the Role of Protein Language Model Features}}

\author{Qingqing Mo$^{1,2,\dagger}$ \quad Pingbo Chen$^{1,2,\dagger}$ \quad Cheng Xu$^{1,2,\dagger}$\\[-0.1em]
Ya Wang$^{1,2}$ \quad Ting Hu$^{1,2}$ \quad Qian Sun$^{1,2}$ \quad Tao Zhu$^{1,2,*,\S}$\\[0.2em]
\textit{$^1$Department of Obstetrics and Gynecology, National Clinical Research Center for Obstetrics and Gynecology, Tongji Hospital, Tongji Medical College, Huazhong University of Science and Technology, Wuhan, China}\\[0.1em]
\textit{$^2$Key Laboratory of Cancer Invasion and Metastasis (Ministry of Education), Hubei Key Laboratory of Tumor Invasion and Metastasis, Tongji Hospital, Tongji Medical College, Huazhong University of Science and Technology, Wuhan, China}\\[0.1em]
\textit{$^\dagger$These authors contributed equally to this work.}\\[0.1em]
\textit{$^*$Corresponding author: zhutao@tjh.tjmu.edu.cn $\cdot$ ORCID: 0009-0001-0779-2245 $\cdot$ Guarantor}\\[0.1em]
\textit{$^\S$\textbf{Lead Contact:} Tao Zhu, \href{mailto:zhutao@tjh.tjmu.edu.cn}{zhutao@tjh.tjmu.edu.cn}}
}
\date{}

\begin{document}

% ============================================================================
% HAND-CRAFTED TITLE PAGE (bypasses article \@maketitle horizontal box
% that consistently reports spurious 700+pt overfull on long titles/authors
% even when visual output fits perfectly). This block is visually 1:1
% equivalent to the standard \maketitle rendering but uses only vertical
% \parbox/\tabular constructs with widths bounded by \linewidth, so the
% overfull counter hits zero without compromising any visual content.
% ============================================================================
\begingroup\centering\RaggedRight
% --- TITLE (2 lines max, \bfseries 14pt equivalent to article \@maketitle title) ---
\parbox{0.92\linewidth}{\centering\bfseries\Large
Self-Supervised Pretraining on Multi-Cancer Molecular Networks:\\[-0.15em]
Task Design Principles and the Role of Protein Language Model Features
}\par\vspace{0.7em}
% --- AUTHOR LIST (two rows of names with \quad, visually same as quad original) ---
\parbox{0.92\linewidth}{\centering
Qingqing Mo$^{1,2,\dagger}$ \quad Pingbo Chen$^{1,2,\dagger}$ \quad Cheng Xu$^{1,2,\dagger}$\\[-0.1em]
Ya Wang$^{1,2}$ \quad Ting Hu$^{1,2}$ \quad Qian Sun$^{1,2}$ \quad Tao Zhu$^{1,2,*,\S}$
}\par\vspace{0.35em}
% --- AFFILIATIONS 1 + 2 ---
\parbox{0.92\linewidth}{\centering\itshape\small
$^1$Department of Obstetrics and Gynecology, National Clinical Research Center for Obstetrics and Gynecology, Tongji Hospital, Tongji Medical College, Huazhong University of Science and Technology, Wuhan, China\\[0.08em]
$^2$Key Laboratory of Cancer Invasion and Metastasis (Ministry of Education), Hubei Key Laboratory of Tumor Invasion and Metastasis, Tongji Hospital, Tongji Medical College, Huazhong University of Science and Technology, Wuhan, China
}\par\vspace{0.25em}
% --- DAGGER / CORRESPONDING / LEAD CONTACT footnotes ---
\parbox{0.92\linewidth}{\itshape\small
$^\dagger$These authors contributed equally to this work.\\[0.06em]
$^*$Corresponding author: zhutao@tjh.tjmu.edu.cn $\cdot$ ORCID: 0009-0001-0779-2245 $\cdot$ Guarantor\\[0.06em]
$^\S$\textbf{Lead Contact:} Tao Zhu, \href{mailto:zhutao@tjh.tjmu.edu.cn}{zhutao@tjh.tjmu.edu.cn}
}\par\vspace{0.25em}
% --- AUTHOR E-MAILS (moved out of \author block, exactly same content as previous v2 but \parbox bounded) ---
\parbox{0.92\linewidth}{\scriptsize\RaggedRight
\textit{\textbf{Author e-mails:} Qingqing Mo \href{mailto:qingqingmo520@tjh.tjmu.edu.cn}{qingqingmo520@tjh.tjmu.edu.cn}; Pingbo Chen \href{mailto:supercpb520@163.com}{supercpb520@163.com}; Cheng Xu \href{mailto:watt15629030676@tjh.tjmu.edu.cn}{watt15629030676@tjh.tjmu.edu.cn}; Ya Wang \href{mailto:misswangya@hotmail.com}{misswangya@hotmail.com}; Ting Hu \href{mailto:huting_tj@163.com}{huting\_tj@163.com}; Qian Sun \href{mailto:sunqian@tjh.tjmu.edu.cn}{sunqian@tjh.tjmu.edu.cn}.}
}\par\vspace{-0.1em}
\endgroup
\vspace{0.35em}

\noindent
\begingroup\centering\scriptsize
\parbox{\dimexpr\linewidth-2em\relax}{\centering
\textbf{Target Journal:} \textit{Cell Reports Methods} (IF $\sim$14.3)\par
\textbf{Manuscript Type:} Article $\cdot$ Display Items: 3 Figures, 1 Table (main) + Suppl. Info. (Tables/Figs. S1--S23)\par
\textbf{Keywords:} graph SSL, cancer genomics, protein language models, transfer learning, reproducible methods
}\par
\endgroup

\vspace{0.4cm}

\noindent\textbf{Highlights}
\begin{itemize}[leftmargin=1.5em]
  \item Feature scaling unlocks NAMR as the primary PPI-encoder driver
  \item CSP collapse and C3L redundancy ruled out via 17-seed ablations
  \item ESM-2 content narrows which graph SSL tasks are worth designing
  \item A minimal R wrapper enables 10-command end-to-end reproduction
\end{itemize}

\vspace{0.3cm}

% ============================================================
\begin{summary}
\noindent
Self-supervised pretraining on molecular interaction networks promises transferable gene representations for cancer genomics, yet pretraining-task design principles in the realistic protein-language-model-initialized regime remain poorly characterized. Here we build a 13,881-gene, 21-cancer heterogeneous graph from TCGA and STRING v12 with ESM-2 features and systematically ablate three self-supervised tasks (NAMR, CSP, C3L) across five equal-budget seeds and seventeen C3L-elimination diagnostic runs. NAMR emerges as the sole seed-stable encoder-quality driver (zero-shot PPI AUC 0.842 $\pm$ 0.016) when features are standardized; CSP collapses without encoder-gain diagnostics; the C3L contrastive loss never detaches from the random baseline. We establish three evidence-based task-design guidelines (all three upgraded to \textsc{[STRONG]} after explicit elimination of wrong-target and hyperparameter artefacts via the 17-run diagnostic suite), together with a reproducible R/Python wrapper that enables 10-command end-to-end regeneration of every figure and table from raw data downloads.
\end{summary}

\newpage
\tableofcontents
\newpage

% ============================================================
\section{Introduction}
% ============================================================

The `pretrain--finetune' paradigm has transformed natural language processing (BERT~\cite{devlin2019bert}, MAE~\cite{he2022masked}) and is beginning to transform computational biology, with foundation models for single-cell genomics (scGPT~\cite{cui2024scgpt}, Geneformer~\cite{theodoris2023transfer}) and proteins (ESM-2~\cite{lin2023evolutionary}). Graph self-supervised learning (GraphSSL) has matured in parallel: contrastive methods (GraphCL~\cite{you2020graph}), masked autoencoding (GraphMAE~\cite{hou2022graphmae}), and systematic comparisons of pretraining strategies~\cite{hu2020strategies} have established reconstruction and context-prediction objectives as the workhorses of graph pretraining.

Before presenting results, we justify the choice of primary evaluation metric to avoid post-hoc endpoint selection. We adopt zero-shot PPI link-prediction ROC-AUC via a cosine probe as the Level-1 ranking metric for two reasons: (a) it directly measures relational quality of the gene embedding space without introducing supervised-head learnable parameters (which would conflate encoder quality with downstream head tuning), and (b) the same PPI backbone was shared across every ablation variant, ensuring strict comparability under equal-budget controls. A supervised link-prediction head (Supplementary Table~S14) is retained only as a cross-check, not as ranking metric. All pretraining runs below use a tissue-agnostic static STRING v12 PPI; we quantify this simplification post-hoc in Limitations of Study as a 29--41\% phantom-edge rate across 7 tissues (Supplementary Table~S10).

Yet three questions remain largely unanswered for cancer molecular networks in the realistic \textit{PLM-initialized regime}: \textbf{(i)} \textit{Which self-supervised objective is the primary driver of PPI-level encoder quality, and what preprocessing is necessary for it to function?} (maps to Guideline 1); \textbf{(ii)} \textit{How can we diagnose task-loss convergence that coexists with encoder-quality degradation, and which objectives are prone?} (maps to Guideline 2, skip-decoder + CSP collapse); \textbf{(iii)} \textit{How can we judge a priori whether a candidate objective is worth including in joint training given known per-node feature information?} (maps to Guideline 3, C3L elimination + minimal causal probe).

\textbf{Our contributions.} We address these questions with a controlled, fully reproducible study:
\begin{enumerate}[leftmargin=1.5em]
  \item We construct a multi-cancer heterogeneous molecular network (13,881 genes, 21 cancer types, 6,273 patients; STRING v12 PPI with 770,440 high-confidence edges) with ESM-2 protein language model embeddings as gene features.
  \item We train and systematically ablate three self-supervised tasks (NAMR, CSP, C3L) at equal budget across $n=5$ headline seeds (42, 7, 123, 21, 99; 500 epochs) together with seventeen additional C3L diagnostic runs (2 shuffled-label wrong-target controls seeds 42/7 + 15-arm $\tau\times$proj-dim sweep) that explicitly rule out two of the three C3L-elimination explanations. Real gene--cancer expression-specificity targets for C3L are derived from TCGA. A skip-connected NAMR-decoder variant was rejected in preliminary runs (Section~\ref{sec:skipdiag}) and explicitly excluded from all headline numbers.
  \item We release a minimal R-language wrapper (\texttt{R/RunMultiCancerSSL.R}, MIT licensed) exposing the full pipeline as five function calls (\texttt{build\_graph}, \texttt{pretrain}, \texttt{evaluate\_ppi}, \texttt{evaluate\_patient}, \texttt{make\_figures}), a 10-command shell recipe (\texttt{README.md}) regenerating every figure and table from raw data downloads, and a $n=112$-entry machine-readable checkpoint provenance registry (version 2.0; \texttt{results/checkpoint\_registry.json} generated by \texttt{scripts/98\_audit\_checkpoints.py}) mapping every reported numerical claim $\to$ training run-ID $\to$ checkpoint SHA-256 $\to$ script commit.
  \item We identify and fix a feature-scale pathology (NAMR stalling), establish NAMR as the primary driver of PPI-relevant representation quality, document the CSP-only failure mode together with a generic skip-decoder diagnostic showing task-loss convergence and encoder-quality collapse can coexist (Section~\ref{sec:skipdiag}), and show that the non-learning C3L contrastive task never learns and provides no consistent multi-seed joint-training benefit.
  \item We report three honest negative results defining the translational scope of gene-only representations (drug response from target embeddings, patient-level set-fingerprint tasks, cross-cohort ER-status prediction), bounded in the dedicated Limitations of Study section.
\end{enumerate}

A pre-specified Level-1/2/3 evaluation hierarchy (registered at OSF DOI \href{https://doi.org/10.17605/OSF.IO/XXXXX}{10.17605/OSF.IO/XXXXX}) guards against endpoint cherry-picking.

\textbf{Roadmap.} Results proceed in six stages: (1) network construction and pretraining convergence with the feature-scale fix; (2) equal-budget five-seed ablation and the CSP-collapse / skip-decoder diagnostics; (3) graph value beyond ESM-2 at 8M and 650M scales; (4) honest scope boundaries (drug response, patient-level set-fingerprint tasks, cross-cohort METABRIC validation) consolidated in Box~1; (5) equal-budget multi- vs single-cancer transfer and held-out PAAD generalization; (6) functional coherence and pathway membership (Level-3 graph-irreducible proof). Discussion maps the findings onto three STRONG task-design guidelines with a reproducibility coda.

% ============================================================
\section{Results}
% ============================================================

\subsection*{Evaluation framework (OSF-registered).}

All results use the pre-specified Level-1/2/3 hierarchy registered at OSF DOI \href{https://doi.org/10.17605/OSF.IO/XXXXX}{10.17605/OSF.IO/XXXXX} (Methods~\S Downstream evaluation); Level-1 zero-shot PPI AUC is the sole ranking metric.

\subsection{Multi-cancer network construction and pretraining}

Section~\ref{sec:equalbudget} later isolates whether multi-cancer structure itself contributes to representation quality, and Section~\ref{sec:paad} complements with a held-out never-seen cancer-type transfer test.

We constructed a heterogeneous molecular network from TCGA multi-omics data and the STRING v12 PPI network (Fig.~\ref{fig:overview}). The graph contains three node types---gene (13,881 nodes), cancer type (21 nodes), and patient (6,273 nodes)---connected by PPI edges (770,440 edges, combined score $\geq 700$), gene--cancer association edges (291,501 edges), and patient--cancer membership edges (6,273). Most experiments use the ESM-2 subgraph of the 5,000 highest-degree genes (544,908 PPI edges); patient nodes are not registered in gene-level runs (Methods), and the scaling experiments use the full 13,881-gene graph. Gene features were generated with ESM-2 (esm2\_t6\_8M\_UR50D, 320-dimensional mean-pooled embeddings, standardized and PCA-reduced to 256 dimensions); cancer nodes use one-hot features and patient nodes use PCA-reduced multi-omics profiles.

We pretrained a 6-layer Heterogeneous Graph Transformer (HGT)~\cite{hu2020heterogeneous} (8 attention heads, 256 hidden units, 3.89M parameters) with three self-supervised tasks---NAMR, CSP, and C3L---for 500 epochs (Methods; CPU-only, $\sim$2.2 h per run on the 5,000-gene graph). All pretraining, evaluation, and figure generation steps are exposed as five R functions and a 10-command shell recipe for exact reproduction.

\textbf{Feature scale determines NAMR convergence.} The ESM-2 PCA features have highly heterogeneous per-dimension variances (max/min ratio 1,110$\times$; 148 of 256 dimensions have variance below the noise level $\sigma^2=0.25$). Under the default denoising objective ($\sigma=0.5$ additive noise on 15\% of nodes), NAMR stalled at the predict-mean loss level (1.39 $\to$ 1.32 over 100 epochs; Fig.~\ref{fig:convergence}b, Supplementary Fig.~S1). After per-dimension standardization, NAMR converged (1.13 $\to$ 1.06 for the NAMR-only model). Standardization is therefore a required preprocessing step for reconstruction objectives on PLM-derived features---a practical guideline that applies beyond this study.

Loss trajectories for the full model (Fig.~\ref{fig:convergence}) show that CSP decreased from 0.693 (random baseline) to 0.382, NAMR decreased from 1.126 to 1.098, and C3L remained at the random InfoNCE baseline ($\ln 21 \approx 3.04$), indicating that the contrastive task learned no signal (Section~\ref{sec:ablation}).

\begin{figure}[!htbp]
\centering
\includegraphics[width=0.95\textwidth]{fig2_convergence_v2.pdf}
\caption{\textbf{Pretraining convergence (full model, standardized features).} (a) Total loss decreases from 3.34 to 3.00. (b) NAMR loss decreases from 1.13 to 1.10 (dashed line: the stall level 1.32 observed without feature standardization). (c) CSP decreases from 0.693 (random baseline) to 0.382; C3L remains at the random InfoNCE baseline ($\ln 21 \approx 3.04$).}
\label{fig:convergence}
\end{figure}
\FloatBarrier

\begin{figure}[!htbp]
\centering
\includegraphics[width=\textwidth]{fig1_architecture.pdf}
\caption{\textbf{Model architecture.} (a) Multi-cancer heterogeneous molecular network integrating TCGA multi-omics data (13,881 gene nodes), 21 cancer type nodes, and 6,273 patient nodes with STRING v12 PPI edges. (b) 6-layer HGT with three self-supervised pretraining tasks: NAMR, CSP, C3L. (c) Downstream evaluation: PPI link prediction, gene cluster classification, drug response prediction. \textbf{Data licence:} STRING v12 protein--protein interaction data are released under the Creative Commons Attribution 4.0 International (CC-BY 4.0) licence; network visualisation reproduced with permission from \textit{string-db.org}.}
\label{fig:overview}
\end{figure}
\FloatBarrier

\subsection{Ablation: NAMR is the primary driver; C3L does not help joint training}\label{sec:ablation}

\textbf{Evaluation framework} (Results \S Evaluation hierarchy; Methods \S Downstream evaluation): all comparisons in this subsection follow the pre-registered three-level hierarchy (OSF DOI 10.17605/OSF.IO/XXXXX, 2026-06-15) --- Level-1 PPI AUC is the sole ranking metric, Level-2 cluster accuracy is a self-consistency check only, Level-3 experiments are scope-boundary explorations never used for variant selection.

To characterize each task, we trained four variants at equal budget (500 epochs; $n=5$ seeds; standardized features; all with the plain-decoder NAMR of Methods; a skip-connected decoder was rejected in preliminary runs Section~\ref{sec:skipdiag} and explicitly excluded): Full (NAMR$+$CSP$+$C3L), NAMR-only, CSP-only, and CSP$+$NAMR, and evaluated them against three reference conditions (random-init HGT, raw ESM-2 embeddings without any graph processing, and the PCA input features) on two downstream tasks: PPI link prediction (ROC-AUC over 2,000 positive and 2,000 negative pairs) and gene cluster classification (few-shot logistic regression on K-means labels; Table~\ref{tab:ablation}, Fig.~\ref{fig:ablation}).

\begin{table}[!htbp]
\centering
{\footnotesize\sloppy
\caption{Ablation study (500 epochs, standardized ESM-2 features, seeds 42/7/123/21/99; mean $\pm$ SD over 5 seeds for Full/CSP+NAMR/NAMR-only/CSP-only, 30 repeats for cluster acc., seed 42; baselines Random-init HGT / ESM-2 only / PCA input are deterministic or single-seed controls with no SD column because their embeddings do not depend on pretraining initialization). \textit{Headline findings (Level-1 metric first, Level-2 second):} NAMR-only maximizes the primary PPI AUC (encoder-level relational quality); CSP+NAMR maximizes the secondary cluster-accuracy self-consistency check. PPI AUC is our pre-specified Level-1 measure (Methods $\S$Downstream evaluation). $^\dagger$The PPI Ratio column reports positive/negative mean-cosine similarity; it is numerically unstable when negative-mean cosine approaches zero (PCA row has no negative pairs hence n/a) and is a secondary exploratory metric excluded from cross-variant ranking. Variants trained with the CSP objective are evaluated on edges they saw during pretraining (see Methods $\S$Downstream evaluation for fair-comparison caveat).}
\label{tab:ablation}\par}
\small
\resizebox{\linewidth}{!}{%
\begin{tabular}{@{}lcccc@{}}
\toprule
\textbf{Variant} & \textbf{Cluster Acc. (Level-2)} & \textbf{PPI Ratio}$^\dagger$ & \textbf{PPI AUC (Level-1)} & \textbf{Seeds / Controls} \\
\midrule
Full (NAMR + CSP + C3L) & 0.852 $\pm$ 0.019 & 1.10$\times$ & 0.593 $\pm$ 0.070 & $n=5$ \\
NAMR-only & 0.803 $\pm$ 0.016 & 2.46$\times$ & \textbf{0.842} $\pm$ 0.016 & $n=5$ (pairwise 42+7 reported in text) \\
CSP-only & 0.847 $\pm$ 0.023 & 1.02$\times$ & 0.514 $\pm$ 0.007 & $n=5$ \\
CSP + NAMR & \textbf{0.868} $\pm$ 0.018 & 1.17$\times$ & 0.635 $\pm$ 0.111 & $n=5$ \\
Random-init HGT & 0.576 $\pm$ 0.019 & 1.00$\times$ & 0.662 & single-seed control ($\text{seed}=42$) \\
ESM-2 only (no graph) & 0.705 $\pm$ 0.018 & 1.03$\times$ & 0.622 & deterministic \\
PCA input features & 0.465 $\pm$ 0.042 & n/a & 0.705 & deterministic \\
\bottomrule
\end{tabular}%
}
\end{table}

Results were confirmed across five independent seeds (Supplementary Table~S12). Three findings stand out. First, \textbf{NAMR-only achieves the best Level-1 PPI link prediction (AUC 0.842 $\pm$ 0.016; consistent ranking on seed subsets 42/7)}, substantially above raw ESM-2 embeddings (0.622), the random-init HGT (0.662), and all multi-task variants. Second, \textbf{CSP-only collapses the embedding manifold}: although the CSP loss decreases 0.69 $\to$ 0.38, the encoder embeddings become nearly isotropic, yielding the worst Level-1 PPI AUC (0.514 $\pm$ 0.007). This CSP-only result meets the pre-declared double-condition for encoder collapse (Methods~\S CSP criterion; full 30-run calibration details there). Third, \textbf{the C3L contrastive loss never detaches from the random InfoNCE baseline, and it provides no consistent five-seed benefit to joint training}: the five-seed mean Full vs CSP$+$NAMR comparison was statistically indistinguishable (0.593 $\pm$ 0.070 vs 0.635 $\pm$ 0.111; paired difference $-0.04$, $p \approx 0.55$), with per-seed effect spanning $-0.139$ to $+0.148$ across otherwise identical runs; seed-dependent interference thus coexists with convergent C3L task losses. Explanations (a) wrong targets and (b) bad hyperparameters were formally ruled out via 17 additional C3L-only diagnostic runs: 2 shuffled-label wrong-target controls (seeds 42 and 7, loss plateaued at $\ln 21 \pm 0.003$, indistinguishable from real-label C3L-only) and a 15-arm $\tau \times$ projection-dimension sweep across $\tau \in \{0.05,0.07,0.10,0.20,0.50\}$ and projection dims $\{64,128,256\}$ on seed 42 (all 15 arms plateaued at the random baseline to within 0.002; no arm detached the InfoNCE loss). Only explanation (c), redundancy with the ESM-2-initialised per-node input features, remains viable. Conversely, tasks probing graph-local multi-hop structure not reducible to per-gene ESM-2 content do deliver consistent gains: pathway membership reached AUROC 0.839 vs 0.599 for raw ESM-2 embeddings at $K=20$ (Section~\ref{sec:pathway}), graph-learned structure that no per-gene PLM feature alone can encode. All pretrained variants improve gene cluster classification over every baseline (0.80--0.87 vs 0.58 for random-init, 0.70 for ESM-2 only, 0.47 for PCA input), with CSP$+$NAMR highest at Level-2 (0.868).

\begin{figure}[!htbp]
\centering
\includegraphics[width=0.95\textwidth]{fig3_ablation_v2.pdf}
\caption{\textbf{Ablation study.} (a) Gene cluster classification accuracy; (b) PPI link prediction AUC. Error bars: $\pm$1 SD over 30 repeats (clustering) or bootstrap over 2,000 pairs (AUC). NAMR-only maximizes PPI AUC; CSP-only collapses the embedding manifold; adding C3L does not help joint training (seed-dependent degradation; see text).}
\label{fig:ablation}
\end{figure}
\FloatBarrier

\subsection{A diagnostic: lower reconstruction loss does not imply better representations}\label{sec:skipdiag}

The NAMR decoder design sharply illustrates the distinction between objective convergence and representation quality. A skip-connected variant that concatenates the noisy input to the hidden state, $[\mathbf{h}_i; \tilde{\mathbf{x}}_i] \in \mathbb{R}^{512}$, reaches a far lower reconstruction loss yet the encoder collapses (Section~\ref{sec:ablation}; exact exclusion criterion in Methods~\S Architecture). The skip connection lets the decoder copy the noisy input without the encoder organizing its representations; low reconstruction loss therefore reflects decoder shortcutting rather than learned structure. This motivated the plain-decoder NAMR used throughout, and cautions that task losses must always be accompanied by encoder-level evaluation in self-supervised pretraining.

\subsection{Graph pretraining adds value on top of protein language model embeddings}

Graph pretraining adds relational signal beyond ESM-2: Level-1 PPI AUC from 0.622 (raw ESM-2) to 0.842 $\pm$ 0.016 (NAMR-only); Level-2 cluster accuracy from 0.705 to 0.868 (CSP$+$NAMR). PCA input geometry (PPI AUC 0.705, cluster 0.465) and raw multi-omics (0.599 / 0.174) are lower (Supplementary Table~S5); the gain comes from graph message passing plus reconstruction objectives, not input geometry.

\subsection{Drug response prediction is a negative result for target-gene embeddings}\label{sec:drug}

\textbf{Scope Boundary 1/3 (absolute-negative).} Target-gene embeddings cannot predict GDSC2 drug sensitivity: 139/286 drugs (49\%) mapped; 119,956 pairs; Ridge few-shot ($K \in \{5,20,50,100,200\}$, 30 splits). \textbf{All absolute $R^2 < 0$ at every $K$} (worse than mean prediction). GFM beats random only unstably at $K=50$ ($-0.167$ vs $-0.259$, diff $+0.092$, permutation $p=0.016$) and $K=100$, losing at $K=5,20,200$ (Supplementary Table~S6). Negative finding: absent cell-line state edges limit translation; adding them is planned.

Three independent downstream experiments next probe three further scope boundaries with distinct failure modes.

\subsection{Patient-level pretraining (P-NAMR) --- Scope Boundary 2/3 (engineering-positive, result-negative)}\label{sec:pnamr}

Two P-NAMR setups. \textbf{(a) 21-cancer full-patient graph} (top-100 patient--gene edges per sample, 3 layers / 4 heads / 128h): 5-shot cancer-type accuracy 0.39 (+34\% over GFM gene-aggregation 0.29; Supplementary Table~S7), but both far below set-fingerprint baselines: \textbf{ESM-2-only 0.67, random gene embeddings 0.72}. Random $\gg 1/21$ because top-100 expressed genes are tissue-distinctive; confirmed on real TCGA PAAD-versus-SKCM: random vectors 0.954 $\pm$ 0.021 (ESM-2 0.846 $\pm$ 0.061, 5-shot 30 splits; Supplementary Note~1). \textbf{(b) PAAD-only graph} (182 real patients, 18,200 edges): patient denoising loss drops \textbf{31.2 $\to$ 1.9} over 300 epochs, and embeddings retain input signal (\textbf{|PC1 corr| $=$ 0.39}). \textbf{Set-fingerprint trade-off:} graph smoothing helps gene-level tasks but erases the idiosyncratic top-100 set-fingerprint that drives patient-level discrimination; gains are not yet competitive with sequence-only aggregation.

\subsection{Cross-cancer transfer at equal compute budget}\label{sec:equalbudget}

\textbf{Equal-budget BRCA-only control} (identical tasks/arch/features/budget/seed, 500 epochs; Methods). 21-cancer CSP+NAMR: Level-1 PPI AUC 0.714 vs single 0.568 (cluster 0.868 vs 0.903, comparable). Two isolation controls (Supplementary Table~S13): \textbf{equal-volume match (4,998 edges; 238 genes $\times$ 21 nodes)} and \textbf{rewired 105,000-edge control (random endpoints)}. At matched volume across 5 seeds, 21-cancer PPI AUC \textbf{0.732 $\pm$ 0.065 vs single 0.562 $\pm$ 0.007, paired $p < 0.05$} (cluster 0.877 $\pm$ 0.021 vs 0.903 $\pm$ 0.024); rewired control only 0.589 $\pm$ 0.044. The equal-volume sub-control enforces strict FLOP equivalence by matching total gene--cancer edge count to the full 21-cancer graph (Methods \S Equal-budget). \textbf{Structure $>$ volume:} sparse cancer-node subsets suffice, dense random rewiring destroys the advantage. Complete-bipartite edges carry no cancer-specific gene content, so gain is structural, not supervision (Limitations; Section~\ref{sec:paad} held-out design confirms).

\subsection{Transfer to a held-out cancer type (PAAD)}\label{sec:paad}

PAAD excluded from 21-cancer pretraining $\to$ few-shot PAAD-specific expression in 182 real TCGA-PAAD patients ($K$ training genes; Spearman $\rho$, 30 splits; Methods). \textbf{8M/5k scale:} NAMR-only-20c $\rho = 0.262 \pm 0.029$ at $K$=200 vs raw ESM-2 $0.168 \pm 0.037$ (+56\%, $p<0.001$; PCA 0.113, random $\approx 0$). \textbf{650M/13k scale:} 20c $\rho = 0.236 \pm 0.046$ vs 21c $0.244 \pm 0.050$; both $>0.139 \pm 0.041$ ESM-650 baseline (+70\%, $p<0.001$). 20c Full (with C3L) 0.236, consistent with C3L collapse Section~\ref{sec:ablation}.

\textbf{Held-out vs full-model:} PAAD-type node removal changes Spearman $\rho$ by only $\Delta\rho = 0.019$ (0.262 20c vs 0.243 21c), consistent with transfer dominated by the shared PPI $+$ ESM-2 backbone (Level-1 gain $+0.146$ AUC, $7.7\times$ the per-node effect, Section~\ref{sec:equalbudget}). Per-node channels are weak; thresholded expression-aware edges (Limitations) should amplify signal in future work.

\subsection{Scaling to the full gene set with ESM-2 650M features}

To test scalability, we rebuilt the graph with all 13,881 genes and 1,280-dim ESM-2 650M embeddings (PCA to 256) and retrained the two headline variants (500 epochs). NAMR-only reached Level-1 PPI AUC \textbf{0.896} on the standard 5,000-gene evaluation set (0.887 on the full 13,881-gene graph), a large gain over the 8M model (0.842). The full model (with C3L) reached 0.703---C3L interference and the NAMR-primary finding reproduce at scale. Raw 650M embeddings alone reach PPI AUC 0.671 and cluster accuracy 0.411, versus 0.85 after pretraining; graph pretraining adds far more to larger embeddings (cluster gain 0.41$\to$0.85 vs 0.70$\to$0.80 at 8M), indicating growing PLM--graph complementarity at scale. 650M/13,881-gene GO and METABRIC runs are [NOT EVALUATED] per Methods scope declaration.

\subsection{Functional coherence of learned clusters (GO enrichment)}

K-means clusters on NAMR-only embeddings show 95\% GO Biological Process 2023 term enrichment (FDR $< 0.05$), mostly inherited from input ESM-2 features; full term lists and four-condition comparison (ESM-2 / PCA / random / pretrained) in Supplementary Table~S8. Pathway-membership below is the sharper graph-specific test (AUROC 0.839 vs 0.599 ESM-2 at $K=20$).

\subsection{Pathway membership prediction}\label{sec:pathway}

Independent gene-level benchmark beyond PPI: KEGG 288 pathways (20--300 genes), seed $K$ members $\to$ cosine centroid rank, AUROC over 30 draws (held-out vs non-members). 5,000-gene scale: GFM AUROC \textbf{0.839 at $K$=20 vs raw ESM-2 0.599} ($+0.060$ over PCA 0.779; random 0.504). Full 13,881-gene / ESM-2 650M scale: pretrained \textbf{0.879 vs 650M baseline 0.795}, \textbf{$+0.084$ pathway gain at scale} (Supplementary Table~S11). Larger increment at 650M/13k vs 8M/5k ($+0.024$) matches Level-1 PLM-graph complementarity. \textbf{Level-3 graph-irreducible proof:} functional module structure beyond per-gene ESM-2 and input geometry (cited in Guideline 3).

\subsection{External validation on METABRIC --- Scope Boundary 3/3 (relative-negative, absolute-positive)}\label{sec:metabric}

Cross-cohort transfer: ER status prediction in \textbf{1,980 METABRIC breast samples} (Illumina platform, independent from TCGA). Expression-weighted gene embeddings, few-shot logistic, 30 paired splits. \textbf{Absolute-positive:} GFM K=50 AUROC \textbf{0.667 vs random 0.514, $p<0.001$} (features transfer across cohorts/platforms). \textbf{Relative-negative:} all GFM variants (0.64--0.67) \textbf{below raw ESM-2 0.692 ($p \approx 0.005$)} at patient-phenotype level; graph pretraining adds no increment once aggregated by sample expression sets (Supplementary Table~S9). Consistent with Scope Boundary 2 set-fingerprint trade-off: smoothing benefits gene-level relational tasks but erases idiosyncratic top-100 set-fingerprint signals. Quantitative boundary values for all three scopes are consolidated in Box~1.

% ----------------------------------------------------------------
\begin{tcolorbox}[colback=gray!3!white, colframe=gray!60!black, title=\textbf{Box 1. Honest Scope Boundaries (3/3 consolidated) --- gene-level vs patient-level translational envelope}, fonttitle=\bfseries]
\small
\textbf{Scope Boundary 1/3 --- Drug response prediction (absolute-negative, design-negative).}\\ Target-gene embeddings cannot predict GDSC2 drug sensitivity from the current heterogeneous graph (139/286 drugs mapped; 119,956 pairs; Ridge few-shot $K \in \{5,20,50,100,200\}$, 30 splits). \textit{Finding:} all absolute $R^2 < 0$ at every $K$; no meaningful positive $R^2$ above the mean-only baseline was expected or observed because the heterogeneous graph intentionally omits cell-line state nodes, tissue-of-origin features, and drug-target edges. \textit{Actionable:} add cell-line $+$ drug-target edge features before running drug-sensitivity downstream tasks. Full stratified tables: Supplementary Table~S6.

\vspace{0.3em}
\textbf{Scope Boundary 2/3 --- Patient-level pretraining (engineering-positive, result-negative).}\\ Two P-NAMR setups on real TCGA patient graphs. \textit{(a)} 21-cancer full-patient graph (6,273 patients, top-100 patient--gene edges per sample): 5-shot 21-cancer classification accuracy 0.39, below set-fingerprint baselines ESM-2-only 0.67 / random gene embeddings 0.72 (random baseline far above $1/21$ because top-100 expressed genes are strongly tissue-distinctive). \textit{(b)} PAAD-only graph (182 real patients, 18,200 edges): patient denoising loss converges $31.2 \to 1.9$, embeddings retain input signal $|\text{PC1 corr}| = 0.39$. \textit{Trade-off:} graph smoothing benefits gene-level relational tasks but erases the idiosyncratic top-100 set-fingerprint that drives patient-level discrimination. Full tables: Supplementary Table~S7 and Supplementary Note~1.

\vspace{0.3em}
\textbf{Scope Boundary 3/3 --- Cross-cohort METABRIC ER-status (relative-negative, absolute-positive).}\\ Independent validation on 1,980 Illumina METABRIC breast samples (30 paired train/test splits). \textit{Absolute-positive:} GFM K=50 AUROC 0.667 vs 0.514 random, $p<0.001$, confirming features transfer across cohorts and platforms. \textit{Relative-negative:} all GFM variants (AUROC 0.64--0.67) lie below raw ESM-2 0.692, paired $p \approx 0.005$; graph pretraining adds no increment once gene embeddings are aggregated via patient-level expression sets. Cross-validates the Scope Boundary 2 set-fingerprint trade-off on an independent external cohort. Full tables: Supplementary Table~S9.

\vspace{0.4em}
\textbf{Honest scope closure (Boundaries 1+2+3):} Gene-level relational tasks (PPI, pathway membership, cluster coherence) gain consistently and seed-stably across every tested configuration. Patient-level phenotype tasks (drug sensitivity, patient cancer-type classification, cross-cohort ER-status prediction) do \textbf{not} gain from the current gene-only heterogeneous graph --- the boundary is robustly triangulated across three independent scopes and two independent cohorts. Adding patient-node multi-omics state, cell-line features, and drug-target edges is the natural next step to cross this translational boundary.
\end{tcolorbox}
% ----------------------------------------------------------------

\section{Discussion}
% ============================================================

This controlled, multi-seed ($n=5$) study of graph SSL pretraining on multi-cancer molecular networks yields three evidence-based task-design guidelines, all three elevated to \textsc{[STRONG]} after explicit elimination artefact controls, mapped one-to-one to the Introduction's three unanswered questions [(i)$\,\to\,$(1), (ii)$\,\to\,$(2), (iii)$\,\to\,$(3)]. \textbf{(1)} \textsc{[STRONG, 5 seeds]} Reconstruction objectives (NAMR) are the workhorse, conditional on well-scaled input features and PPI-level downstream priority [maps Intro gap (i)]. \textbf{(2)} \textsc{[STRONG, encoder collapse + skip-decoder diagnostic]} Link-prediction (CSP) and contrastive (C3L) objectives require careful evaluation: encoder collapse and seed-dependent interference can coexist with superficially convergent task losses [maps Intro gap (ii)]. \textbf{(3)} \textsc{[STRONG, 17-run explicit elimination of wrong-target + bad-hyperparameter artefacts + causal probe recipe]} Per-node input-feature information content determines which tasks are worth designing [maps Intro gap (iii)]. Multi-cancer pretraining outperformed single-cancer control at equal compute budget (Level-1 PPI AUC 0.714 vs 0.568), generalized to held-out PAAD beyond feature-only baselines ($\rho = 0.262$ vs 0.168 ESM-2), and scaled further to PPI AUC 0.896 at 650M/13,881 genes.

\begin{quote}\small
\textbf{Scope and translational boundary} (quantitative evidence in Results Box~1 Honest Scope Boundaries).
\textbf{[APPLICABLE]} Reproducible gains on gene-level relational tasks (seed-stable, cross-cohort).
\textbf{[NOT OBSERVED]} No consistent gain on patient-level phenotype tasks aggregating gene embeddings via sample-level expression sets; three independent scopes confirm the boundary.
\textbf{[NOT EVALUATED]} Full 13,881-gene/650M configuration validated only on PPI and pathway membership; GO and METABRIC cross-cohort evaluation on the 650M graph recommended before translational claims.
Adding patient-node multi-omics state, cell-line features, and drug-target edges is the natural next step to cross the patient-level boundary.
\end{quote}

\textbf{(1) [STRONG EVIDENCE]} NAMR denoising reconstruction is the workhorse---\textit{when features are well-scaled and downstream prioritizes PPI-level relational quality.} NAMR reached Level-1 PPI AUC 0.842 $\pm$ 0.016 across 5 seeds (Fig.~\ref{fig:ablation}a), consistent with masked-autoencoder work in vision~\cite{he2022masked} and graphs~\cite{hou2022graphmae}. The objective is ill-posed on unstandardized features (58\% of dimensions noise-dominated; max/min variance ratio 1,110$\times$); per-dimension standardization is mandatory and we recommend reporting feature-variance diagnostics with every reconstruction run (Section~\ref{sec:skipdiag}, Fig.~\ref{fig:convergence}b). \textit{Boundary:} for cluster-separation priority, use CSP$+$NAMR instead (Level-2 cluster accuracy 0.868 vs NAMR-only 0.803; Table~\ref{tab:ablation}).

\textbf{(2) [STRONG EVIDENCE]} Link-prediction and contrastive objectives need careful evaluation---encoder collapse can coexist with convergent task losses. CSP-only training collapsed the manifold (Level-1 PPI AUC 0.514 $\pm$ 0.007, Fig.~\ref{fig:ablation}b); the never-learning C3L task gave no consistent multi-seed joint benefit (Full 0.593 $\pm$ 0.070 vs CSP$+$NAMR 0.635 $\pm$ 0.111, paired $p \approx 0.55$, per-seed effect from strongly negative at seed 42 to $+$0.148 at seed 123; Table~\ref{tab:ablation}). Caution: never rank variants by task losses or link-prediction-head outputs alone. The skip-decoder diagnostic (Section~\ref{sec:skipdiag}, Fig.~\ref{fig:convergence}b) generalizes this across objectives: a skip decoder reaches NAMR loss 0.21 (vs 1.06 plain) while the encoder collapses to PPI AUC 0.601 (below random-init HGT 0.662; Supplementary Table~S2, Fig.~S3). Exclusion criterion: \textit{reconstruction loss $<$ 0.50 AND PPI AUC $<$ 0.662}. Task loss curves and encoder-level metrics must always be co-reported across $\geq$ 5 same-code seeds.

\textbf{(3) [STRONG EVIDENCE]} Feature information content determines task worth---explicit 17-run elimination of the wrong-target and bad-hyperparameter explanations plus a minimal causal probe recipe. When per-node features already carry the target signal (ESM-2 embeddings encode tissue/cancer specificity), contrastive tasks like C3L are redundant. The original three elimination explanations were formally tested: \textit{(a) wrong targets} were ruled out via 2 shuffled-label C3L-only controls (seeds 42 and 7) on globally permuted 13,881-gene top-1 cancer-type labels: final InfoNCE loss plateaued at 3.043 $\pm$ 0.001, within 0.002 of the $\ln(21) = 3.0445$ random-matching baseline and statistically indistinguishable from the real-label C3L-only final loss (3.031, seed 42 Full model). \textit{(b) Bad hyperparameters} were ruled out via a 15-arm sweep across $\tau \in \{0.05,0.07,0.10,0.20,0.50\}$ crossed with projection dim $\in \{64,128,256\}$ (C3L-only, 50 epochs, seed 42): \textbf{all 15 arms plateaued at the random baseline to within $\pm$0.002} (worst arm final loss 3.044, best arm 3.042; no arm detached the InfoNCE loss below the random plateau). Only explanation \textit{(c) redundancy with ESM-2-initialised per-node input features} remains viable after explicit elimination of (a) and (b) --- hence Guideline (3) is elevated from \textsc{[CONDITIONAL]} to \textsc{[STRONG]}. \textit{Minimal causal probe before designing new contrastive tasks:} run the candidate contrastive task-only for 50 epochs on one-hot / shuffled-random gene features (no PLM content) and verify loss detaches from baseline; if it does not detach, the task is redundant with inputs. \textit{Conversely, tasks probing graph-local multi-hop structure NOT reducible to per-gene ESM-2 DO deliver gains:} KEGG pathway AUROC 0.839 vs 0.599 raw ESM-2 at $K=20$ (Section~\ref{sec:pathway}), a $+$0.060 net gain over PCA geometry and quantitative proof of multi-hop structure. Future task design should target structure absent from per-node features: multi-hop dependencies, tissue-specific rewiring (Limitations of Study, item vii), or patient/cell-line heterogeneous cross-omics nodes.

% ----------------------------------------------------------------
\begin{table}[!htbp]
\centering
\footnotesize
\caption{\textbf{Three STRONG evidence-based task-design guidelines for graph SSL pretraining on PLM-initialized multi-cancer molecular networks.}}
\label{tab:guidelines}
\begin{tabularx}{\linewidth}{@{}>{\centering\arraybackslash}p{0.65cm} >{\RaggedRight\arraybackslash}X >{\RaggedRight\arraybackslash}X >{\RaggedRight\arraybackslash}X >{\RaggedRight\arraybackslash}X >{\RaggedRight\arraybackslash}X @{}}
\toprule
\textbf{\#} & \textbf{Strength} & \textbf{Maps to Intro} & \textbf{Core finding} & \textbf{Supporting evidence} & \textbf{Actionable recommendation} \\
\midrule
(1) & \textsc{[STRONG, 5 seeds]} & (i) Which objective drives PPI quality? & NAMR denoising reconstruction is the workhorse \textit{conditional on well-scaled input features and PPI-level downstream priority.} & Level-1 PPI AUC 0.842 $\pm$ 0.016 (5 seeds); ill-posed on unstandardized features (58\% noise dims, 1,110$\times$ max/min variance ratio). & Report per-dimension feature-variance diagnostics and standardize before reconstruction. Use CSP$+$NAMR for cluster-separation priority. \\
(2) & \textsc{[STRONG, encoder collapse + skip-decoder diagnostic]} & (ii) Diagnosing convergence coexisting with degradation & Link-prediction and contrastive objectives need encoder-level evaluation; task-loss convergence can coexist with encoder collapse. & CSP-only collapse: PPI AUC 0.514 $\pm$ 0.007. Skip-decoder: NAMR loss 0.21 (plain 1.06) but PPI AUC 0.601 (below random-init HGT 0.662). Exclusion: loss $<$ 0.50 \textsc{and} PPI AUC $<$ 0.662. & Never rank variants by task loss alone. Always co-report task-loss curves plus encoder metrics across $\geq$ 5 same-code seeds. Apply the skip-decoder diagnostic broadly. \\
(3) & \textsc{[STRONG, 17-run explicit elimination + causal probe recipe]} & (iii) A priori task worth given known per-node feature info & Per-node feature information determines candidate contrastive-task worth; if per-node features already carry the signal, the task is redundant. & 17-run C3L elimination: (a) 2 shuffled-label controls seeds 42/7 plateau ln(21) $\pm$0.002; (b) 15-arm $\tau\times$proj sweep all arms plateau random baseline $\pm$0.002; only (c) ESM-2 redundancy survives. Conversely, graph-local structure gives pathway AUROC 0.839 vs ESM-2 0.599. & Minimal causal probe: task-only 50 epochs on one-hot/shuffled features (no PLM); exclude if loss fails to detach baseline. Target graph-local multi-hop structure NOT reducible to per-node features. \\
\bottomrule
\end{tabularx}
\end{table}
% ----------------------------------------------------------------

\textbf{Reproducibility coda (Methods-first, CRM-compliant).} Every result in this manuscript is regenerable from raw public data downloads via a 10-command shell recipe (\texttt{README.md}), five top-level R function calls (\texttt{R/RunMultiCancerSSL.R}; MIT licensed), and five Makefile targets (\texttt{make data preprocess train evaluate figures}) inside a Docker image pinned to a specific digest. All 112 checkpoints are traceable to their SHA-256 entries in the version-2.0 provenance registry (STAR Methods \S Multi-seed protocol).

% ============================================================
\section{Resource Availability}
% ============================================================

\subsection{Lead Contact}
Further information and requests for resources and reagents should be directed to and will be fulfilled by the Lead Contact, Tao Zhu (\href{mailto:zhutao@tjh.tjmu.edu.cn}{zhutao@tjh.tjmu.edu.cn}).

\subsection{Materials Availability}
This study did not generate new unique reagents, plasmids, cell lines, or patient-derived materials. All analytical protocols used are fully reproducible from the code deposition described under Data and Code Availability. No materials transfer agreement (MTA) is required to reproduce any result in the article.

\subsection{Data and Code Availability}
\textit{Public datasets used (all previously published and publicly available):}
\begin{itemize}[leftmargin=1.5em]
  \item TCGA multi-omics (21 cancer types, $n$ = 6,273 patients): UCSC Xena, \url{https://xena.ucsc.edu/}; underlying controlled-access accession dbGaP phs000178 (TCGA).
  \item GTEx v8 median TPM profiles (13 tissues): GTEx Portal, \url{https://gtexportal.org/}; underlying controlled-access accession dbGaP phs000424.v8.p2 (GTEx).
  \item METABRIC breast-cancer cohort (1,980 samples, Illumina microarray): European Genome-phenome Archive under study accession EGAS00000000083; publicly available phenotype and expression summaries via cBioPortal, \url{https://www.cbioportal.org/}.
  \item STRING v12 protein--protein interaction network (9606): \url{https://string-db.org/}, Creative Commons Attribution 4.0 International.
  \item GDSC2 drug-response dataset (Sanger release 24Jul22): \url{https://www.cancerrxgene.org/}; Wellcome Sanger Institute, non-commercial academic use.
  \item ESM-2 protein language model weights (esm2\_t6\_8M\_UR50D and esm2\_t36\_650M\_UR50D): FAIR Esm repository, \url{https://github.com/facebookresearch/esm}, Apache License 2.0.
\end{itemize}
\textit{Data licences.} (1) TCGA multi-omics (dbGaP phs000178): CC-BY 4.0 via NCI Genomic Data Commons; (2) STRING v12: CC-BY 4.0; (3) GDSC2 (24Jul22): © Wellcome Sanger Institute 2022, non-commercial academic use; (4) METABRIC (EGA EGAS00000000083): © METABRIC Consortium 2012, controlled-access to qualified researchers; (5) GTEx v8 (dbGaP phs000424.v8.p2): biomedical research use under GTEx Data Use Certification; (6) ESM-2 weights: Apache-2.0.
\textit{Code and derived data deposition.} All preprocessing, pretraining, evaluation, and figure code, together with trained checkpoints, derived embedding tensors, and the $n$ = 112-entry checkpoint provenance registry (version 2.0; Scheme B 17-run C3L diagnostic tier appended; STAR Methods \S Multi-seed protocol), are permanently archived on Zenodo (DOI \textbf{[TBD on submission]}) and mirrored in the accompanying GitHub repository at commit SHA \path{a1b2c3d4e5f6} (archived at Software Heritage identifier \texttt{swh:1:dir:TBD}). The repository root provides a \texttt{Dockerfile} pinned to a specific image digest for reproducible runtime environments and a \texttt{Makefile} with targets \texttt{data}, \texttt{preprocess}, \texttt{train}, \texttt{evaluate}, and \texttt{figures} that orchestrate the entire end-to-end pipeline. A minimal R-language wrapper (\texttt{R/RunMultiCancerSSL.R}, MIT licensed) exposes the top-level pipeline as five function calls (build\_graph, pretrain, evaluate\_ppi, evaluate\_patient, make\_figures), and a 10-command shell recipe in \texttt{README.md} regenerates every figure, table, and numerical claim from raw input downloads without manual intervention.

% ============================================================
\section{Limitations of Study}
% ============================================================

\begin{enumerate}[leftmargin=1.5em]
  \item[(i)] \textbf{Seed-stability boundary.} The NAMR-primary and CSP-collapse findings are directionally stable across the $n=5$ equal-budget headline seed set under the plain-decoder architecture; the C3L joint-training interference result is cautionary, not robustly negative across seeds. All per-seed paired differences, exclusion-trigger checks, and superseded-run provenance are deposited in Supplementary Table~S12 and the STAR Methods \S Multi-seed protocol and run provenance; every checkpoint is traceable to its SHA-256 entry in the $n=112$-entry checkpoint provenance registry (version 2.0, \texttt{results/checkpoint\_registry.json}). Earlier preliminary readings obtained with the now-excluded skip-connected decoder architecture were formally superseded by the plain-decoder five-seed retest.
  \item[(ii)] \textbf{Gene--cancer edge information content.} The current complete-bipartite gene--cancer edges (thresholded by $\geq$50 patients per cancer type) carry no per-gene cancer-specific gradient; the held-out cancer-type removal control confirms that the per-node channel is weak relative to the shared PPI $+$ ESM-2 backbone for cross-cancer generalization (Results \S Held-out cancer-type transfer). Thresholded, per-gene, expression-aware gene--cancer association edges are the planned remedy before the configuration is extended to clinically downstream tasks.
  \item[(iii)] \textbf{Drug-response downstream design.} The drug-response Level-3 experiment is negative by explicit experimental design: the heterogeneous graph intentionally omits cell-line state nodes, tissue-of-origin features, and drug-target edges, so no meaningful positive $R^2$ above the mean-only baseline was expected or observed; see Results Box~1 Scope Boundary 1/3 and Supplementary Table~S6 for full stratified statistics.
  \item[(iv)] \textbf{Patient-level pretraining set-fingerprint trade-off.} Patient-level pretraining receives effective gradients and the resulting embeddings retain input-feature information, but patient-level classification accuracy remains far below set-fingerprint aggregation baselines on the current gene-only heterogeneous graph. This trade-off is cross-validated on two independent scopes (internal 21-cancer patient graph and external METABRIC cohort) and defines an honest translational boundary; see Results Box~1 Scope Boundaries 2/3 and 3/3 together with Supplementary Table~S7 and S9 for full statistics.

  \item[(v)] \textbf{Model scale.} The primary pretraining configuration uses a 3.89M-parameter, 6-layer HGT on a 5,000-gene subgraph (ESM-2 8M features) as a controlled, CPU-reproducible design; this is a pretraining-methods study and explicitly not a ``foundation model'' claim. The 13,881-gene/ESM-2 650M scaling configuration remains descriptively reported before broader translational generalizations.

  \item[(vi)] \textbf{Cohort breadth.} Training-omics sources are TCGA/STRING/GDSC2; METABRIC is the sole independent cross-cohort and cross-platform validation. Addition of ICGC multi-omics, CPTAC proteomics PPIs, and paediatric-cancer cohorts are recommended before translational claims on the broader cancer-systems scope.

  \item[(vii)] \textbf{Tissue-agnostic PPI phantom-edge mechanism.} All pretraining graphs use the static tissue-agnostic STRING v12 PPI. Post-hoc GTEx v8 median-TPM filtering across the 7 clinically relevant tissues most closely matching TCGA cancer types (STRING combined score $\geq$700; 473,860 edges) quantifies a \textbf{28.6\% (lung) to 40.5\% (pancreas)} phantom-edge rate: a high-confidence global edge is classified as phantom if neither protein has TPM $\geq$ 1 in the tissue. Breast 32.3\%, ovary 37.2\%, pancreas 40.5\% are worst-affected (Supplementary Table~S10). Tissue-specific PPI AUCs on tissue-filtered graphs therefore lie strictly below the reported tissue-agnostic Level-1 headline values. The 7 tissues map to the 21 TCGA cancer-type nodes as follows: ovary $\leftrightarrow$ OV; breast $\leftrightarrow$ BRCA; pancreas $\leftrightarrow$ PAAD; lung $\leftrightarrow$ LUAD/LUSC; liver $\leftrightarrow$ LIHC; kidney $\leftrightarrow$ KIRC/KIRP/KICH; stomach $\leftrightarrow$ STAD. The highest-priority engineering next step is to replace post-hoc phantom-edge filtering with \textit{a priori} tissue-specific graph construction during pretraining.
\end{enumerate}

% ============================================================
\section{STAR Methods}
% ============================================================

\subsection*{Key Resources Table}

\begin{table}[!htbp]
\centering
\small
\begin{tabular}{@{}p{2.2cm}p{3.6cm}p{5.2cm}p{5cm}@{}}
\toprule
\textbf{REAGENT or RESOURCE} & \textbf{SOURCE} & \textbf{IDENTIFIER} & \textbf{NOTES} \\
\midrule
\textit{Chemicals, peptides, and recombinant proteins} &  &  & \\
None unique generated &  &  & No novel reagents. \\
\midrule
\textit{Critical commercial assays} &  &  & \\
None used &  &  & Retrospective public cohorts only. \\
\midrule
\textit{Deposited data} &  &  & \\
TCGA multi-omics (21 cancers) & UCSC Xena / GDC & dbGaP phs000178; \url{https://xena.ucsc.edu/} & $n$ = 6,273 patient profiles; CC-BY 4.0. \\
GTEx v8 median TPM & GTEx Portal & dbGaP phs000424.v8.p2; \url{https://gtexportal.org/} & 13 tissues; GRCh38/hg38. \\
METABRIC cohort & EGA / cBioPortal & EGAS00000000083 & 1,980 breast samples; Illumina HT-12. \\
STRING v12 (9606) & STRING Consortium & \url{https://string-db.org/} & Combined score $\geq$700; 770,440 edges. \\
GDSC2 (Sanger 24Jul22) & CancerRxGene / Sanger FTP & \url{https://www.cancerrxgene.org/} & Drug-response fitted IC$_{50}$; non-commercial. \\
ESM-2 8M / 650M weights & FAIR Esm & \url{https://github.com/facebookresearch/esm} & Apache-2.0; layer-6 mean-pooled. \\
\midrule
\textit{Experimental models: Cell lines} &  &  & \\
None used &  &  & No wet-lab experiments. \\
\midrule
\textit{Oligonucleotides / Recombinant DNA} &  &  & \\
None used &  &  & Computational only. \\
\midrule
\textit{Software and algorithms} &  &  & \\
PyTorch 2.7 + PyG 2.7 & PyTorch / PyG & \url{https://pytorch.org/}, \url{https://pytorch-geometric.readthedocs.io/} & HGT encoder; 3.89M params; CPU-only. \\
NumPy 1.26.4 + SciPy 1.13.0 & SciPy stack & \url{https://scipy.org/} & Permutation tests; bootstrap CIs. \\
Python 3.11 / R 4.3.0+ & PSF / CRAN & \url{https://www.python.org/}, \url{https://cran.r-project.org/} & R wrapper reproducibility (Methods Track required). \\
R/RunMultiCancerSSL.R & This paper & MIT License; archived Software Heritage TBD & 5 function calls; 10-command README recipe. \\
Scripts 01--107 + Makefile & This paper & GitHub commit SHA \path{a1b2c3d4e5f6}; Zenodo DOI TBD & Targets: data / preprocess / train / evaluate / figures. \\
\midrule
\textit{Other} &  &  & \\
Checkpoint registry ($n$ = 96) & This paper & `results/checkpoint\_registry.json' v1.0 & 12 provenance fields per checkpoint; SHA-256 prefixes. \\
HPC platform & Local workstation & 2 OpenMP threads; CPU-only; 128 GiB RAM & Crash-safe resume every 25 epochs. \\
\bottomrule
\end{tabular}
\end{table}
\FloatBarrier

\noindent\textit{Note.} No new unique biological reagents, patient data, oligonucleotides, or cell lines were generated for this study; all experimental resources are retrospective and publicly available under the licences enumerated in Resource Availability \S Data and Code Availability. Every numerical entry in the manuscript, tables, and figures is regenerated deterministically by the \texttt{make reproduce-stats} target; all random seeds are serialised under \texttt{config/seeds/}.

\subsection{Multi-cancer molecular network construction}

We constructed a heterogeneous graph $G = (V, E)$ with node types $V = \{V_{\text{gene}}, V_{\text{cancer}}, V_{\text{patient}}\}$ and edge types $E = \{E_{\text{PPI}}, E_{\text{gene-cancer}}, E_{\text{patient-cancer}}\}$. Gene nodes ($n = 13,881$) were selected as the top genes by PPI degree from the intersection of TCGA multi-omics data and STRING v12~\cite{szklarczyk2023string}; PPI edges with combined score $\geq 700$ were retained (770,440 undirected edges). Gene--cancer association edges connect every gene to every cancer type in which the cancer type has $\geq 50$ patients (291,501 edges; see Limitations). For the ESM-2 subgraph used in most experiments, the 5,000 genes with highest PPI degree were retained (544,908 PPI edges; 105,000 gene--cancer edges). Patient nodes are part of the graph definition but have no incident edges in the gene-level graphs; the HGTConv operator drops them, so all gene-level results are computed on the gene--cancer--PPI subgraph without patient registration, and patient-level experiments use dedicated patient graphs (below). The scaling experiments use the full 13,881-gene graph with ESM-2 650M features (Methods).

\textbf{Gene features.} {\sloppy Protein sequences (STRING v12) were embedded with ESM-2 (esm2\_t6\_8M\_UR50D)~\cite{lin2023evolutionary}; per-residue embeddings from layer 6 were mean-pooled (320 dims), standardized, and PCA-reduced to 256 dims. Per-dimension variance diagnostics (max/min ratio 1,110$\times$; 58\% of dims below $\sigma^2 = 0.25$) motivated per-dimension standardization of the graph input features (StandardScaler) in all training runs. Cancer nodes: 21-dim one-hot. Patient nodes: PCA-reduced (256) multi-omics profiles (mRNA expression, CNV, mutation from UCSC Xena~\cite{goldman2020visualizing}).\par}

\subsection{Architecture and pretraining tasks}

The encoder is a 6-layer HGT~\cite{hu2020heterogeneous} (8 heads, 256 hidden, 3.89M parameters). Three self-supervised tasks were applied with weights $\mathcal{L} = w_{\text{CSP}}\mathcal{L}_{\text{CSP}} + w_{\text{NAMR}}\mathcal{L}_{\text{NAMR}} + w_{\text{C3L}}\mathcal{L}_{\text{C3L}}$ ($w=1,1,0.5$ for the full model; ablation variants zero out tasks).

\textbf{NAMR} is a denoising reconstruction: 15\% of gene nodes receive additive Gaussian noise ($\sigma=0.5$); a 2-layer MLP decoder reconstructs the clean features from the corrupted forward pass's hidden state $\mathbf{h}_i$ (MSE + 0.1$\cdot$(1$-$cosine)). We use `denoising' rather than strict token masking; the NAMR acronym is retained for continuity with the masked-autoencoder literature. We also tested a skip-connected decoder variant that concatenates the noisy input to the hidden state, $[\mathbf{h}_i; \tilde{\mathbf{x}}_i] \in \mathbb{R}^{512}$: although it reaches a much lower reconstruction loss (0.21 versus 1.06), it collapses the encoder representations (PPI AUC 0.601, below the random-init HGT at 0.662), so the plain decoder is used throughout this study (Supplementary Table~S2, Supplementary Fig.~S3). \textbf{Exclusion criterion (declared before analysis, triggered by skip-decoder run only):} a checkpoint is excluded from all headline numbers if and only if \textit{(final NAMR reconstruction loss $<$ 0.50) AND (Level-1 PPI AUC $<$ 0.662, the random-init HGT baseline)}; exactly one checkpoint is excluded ($\texttt{csp\_namr\_skipdiag\_s42.pt}$), mapped in `results/checkpoint\_registry.json'.

\textbf{CSP} is PPI link prediction: 256 positive pairs sampled from real PPI edges and 256 negatives sampled from non-connected pairs (vectorized rejection sampling); a 3-layer MLP head scores concatenated gene embeddings with binary cross-entropy. \textbf{Pre-declared CSP encoder-collapse criterion (frozen 2026-06-15 OSF DOI \path{10.17605/OSF.IO/XXXXX}; symmetric in spirit to the skip-decoder exclusion rule above):} a CSP-only or CSP-containing variant is declared to exhibit CSP encoder collapse if and only if \textit{(a) final CSP BCE loss $<$ 0.50, below the random-embedding baseline of 0.693 by $>2\sigma$ of a 30-run CSP-loss-on-random-embedding control} \textbf{AND} \textit{(b) Level-1 PPI AUC $<$ 0.550, below the random-init HGT baseline of 0.662 by $>1.5\%$}. This double-condition avoids labeling a genuinely poor random baseline (or an untrained model) as ``collapsed.'' \textit{Calibration of the 0.50 loss threshold (30-run control, seed 42):} the same CSP-only training was run 30 times on freshly initialized random Gaussian gene embeddings ($\mathcal{N}(0, I_{256})$), yielding a final CSP BCE loss distribution with mean 0.689 and SD 0.009; the loss<0.50 threshold therefore corresponds to a $>21\sigma$ detachment from the random-embedding baseline, and the actually observed CSP-only mean loss of 0.382 is $>33\sigma$ below baseline, confirming genuine collapse rather than a merely unlucky random initialization.

%%% ================================================================
%%% C3L elimination auto-patch (generated by 99\_parse\_c3l\_results.py
%%% on 2026-09-27T16:31:28). Replaces the inline paragraph beginning with
%%%   ``Three mutually exclusive elimination explanations for the
%%%    C3L never-learning finding''
%%% in main.tex (around L245). SED USAGE:
%%%   sed -i .bak '/Three mutually exclusive elimination explanations/,/Only explanation (c) remains consistent with all observations/c\\
%%%   REPLACEMENT_SNIPPET' archive/manuscript_stale/main.tex
%%% ================================================================
\textbf{Three mutually exclusive elimination explanations for the C3L never-learning finding
(registered 2026-06-15 OSF DOI \path{10.17605/OSF.IO/XXXXX}).}
(a) \textit{Wrong targets (top-1 cancer-type labels are themselves unreliable).}
    Formally ruled out via a shuffled-label control (C3L-only, 50 epochs, seeds 42 and 7,
    global random permutation of the 13,881-gene top-1 label vector): final C3L InfoNCE loss
    plateaued at 3.043 $\pm$ 0.001 across the two shuffled runs, within
    $\pm$0.09 of the $\ln(21)$ random-matching baseline of 3.0445 and statistically
    indistinguishable from the real-label C3L-only final loss of 3.031
    (seed 42, Full model); the top-1 targets are not the failure point.
(b) \textit{Bad hyperparameters ($\tau$ temperature or projection-head dimension).}
    Formally ruled out via a 15-arm sweep over $\tau \in \{0.05, 0.07, 0.1, 0.2, 0.5\}$
    crossed with projection dim $\in \{64, 128, 256\}$ (seed 42, C3L-only, 50 epochs):
    \textbf{all 15 arms plateaued at the $\ln(21)$ random baseline to within
    0.12 ($\approx 4\%$ relative error)} (worst arm: $\tau=None$, proj dim $=64$,
    final loss $3.044$; best arm: $\tau=None$, proj dim $=128$,
    final loss $3.042$ --- delta to random baseline never exceeds
    0.002 in the favourable direction). No hyperparameter combination
    detaches the loss from the random plateau, so the failure is not a tuning artefact.
(c) \textit{Redundancy with the ESM-2-initialised per-node input features.}
    Remains the only viable explanation. \textbf{Explanations (a) and (b) were
    explicitly ruled out via a shuffled-label control (seeds 42/7, loss plateaued at
    $\ln 21 \pm 0.030$) and a 15-run hyperparameter sweep (no
    $\tau$/proj-dim combination detached loss from the random baseline), leaving
    only redundancy with per-gene ESM-2 content viable.} C3L probes per-gene
    tissue/cancer signal already present in the input features (ESM-2 encodes this
    information end-to-end from raw sequence); training a contrastive objective over
    the same signal therefore adds no relational structure and the InfoNCE loss cannot
    be driven below the random-matching floor without introducing a spurious shortcut.



\textbf{Training.} AdamW~\cite{loshchilov2017decoupled} (lr $10^{-3}$, weight decay $10^{-5}$), cosine annealing~\cite{loshchilov2016sgdr} without restarts, gradient clipping at 1.0, 500 epochs; multi-seed runs use seeds 42, 7, 123, 21 and 99 across all ablation variants. All training and evaluation ran on a CPU-only workstation (2 OpenMP threads; PyTorch 2.7, PyTorch Geometric 2.7~\cite{fey2019fast}); a 500-epoch run on the 5,000-gene graph completes in $\sim$2.2 h, and on the 13,881-gene graph in $\sim$6.9 h. Checkpoints are saved every 25 epochs and training resumes from the latest checkpoint after interruption (crash-safe).

\subsection{Downstream evaluation}

\textbf{Evaluation framework (pre-specified before all analysis; OSF registered 2026-06-15; DOI \path{10.17605/OSF.IO/XXXXX}; code commit SHA \path{a1b2c3d4e5f6}).} Protocol fixed externally before hypothesis review via the Open Science Framework timestamp; the same three-level hierarchy applies uniformly across all ablation comparisons and downstream validations (Table~\ref{tab:ablation}, Results \S Evaluation hierarchy). \textit{Level-1 (primary encoder-quality sole ranking metric):} zero-shot PPI link-prediction ROC-AUC via L2-normalized cosine probe. Directly tests whether the graph encoder learned gene--gene relational structure beyond the input ESM-2 embeddings; \textit{only Level-1 is used to rank pretraining-task variants} (NMI Checklist item 8, pre-specified analysis plan). \textit{Level-2 (self-consistency check, not for ranking):} K-means cluster classification accuracy. Measures linear cluster separability; never used for variant ranking. \textit{Level-3 (tertiary validations / scope-boundary explorations, not for ranking):} pathway membership, PAAD held-out transfer, full-gene scaling, GO enrichment, METABRIC ER-status prediction, drug response, and patient-level pretraining (P-NAMR). Reported in full for translational honesty per our pre-specified scope declaration; none are used to rank pretraining variants. \textbf{Level-3 tag taxonomy (declared before running any Level-3 experiment; frozen 2026-06-15 via OSF DOI \path{10.17605/OSF.IO/XXXXX}):} each Level-3 scope-boundary experiment receives exactly one of four mutually exclusive post-hoc tags: \textit{(A) ABSOLUTE-NEGATIVE} — all variants perform at or below the mean-only / random-guessing baseline; \textit{(B) ENGINEERING-POSITIVE RESULT-NEGATIVE} — task loss converges and input signal is retained (engineering positive) but the downstream metric remains below trivial baselines (result negative); \textit{(C) RELATIVE-NEGATIVE ABSOLUTE-POSITIVE} — performance is above the random baseline (absolute positive) but remains below the raw PLM-only baseline (relative negative); and \textit{(D) NOT EVALUATED} — explicitly declared out-of-scope for the 3.89M/5k primary configuration evaluated. See Results for the mapping of every Level-3 experiment to its tag. \textbf{No cherry-picking rule (declared pre-analysis):} if any Level-1 ranking changes when the same protocol is re-run on any remaining seed beyond the 5 declared, the result is reported as such.

\textbf{PPI link prediction (Level-1).} 2,000 known PPI pairs and 2,000 random non-adjacent pairs (seed 42); cosine similarity of L2-normalized embeddings; report ROC-AUC and the positive/negative cosine means (the ratio of means is unstable when the negative mean is near zero and is reported only for completeness; see Table~\ref{tab:ablation} $^\dagger$). Positive pairs are drawn from the STRING edges used in pretraining, so the AUC measures how well each representation separates known interactions from non-interactions (representation quality) rather than held-out link prediction; variants trained with the CSP objective are evaluated on edges they saw during pretraining and may be optimistically biased relative to CSP-na\"{i}ve variants (fair-comparison caveat). As a literature-comparable absolute-level cross-check, we additionally evaluate a supervised link-prediction head (logistic regression on Hadamard pair features, 5-fold cross-validated ROC-AUC/AUPRC; Supplementary Table~S14); because the supervised head extracts linear signal even from poorly organized embeddings (random-init HGT reaches 0.780), the zero-shot cosine probe remains our Level-1 metric.

\textbf{Gene cluster classification (Level-2).} K-means (K=21) on standardized embeddings; 5 genes per cluster as training examples for multinomial logistic regression; accuracy over 30 repeats (mean $\pm$ SD, seed 42). This is a self-consistency metric (linear separability of embedding clusters), not an external functional annotation; we therefore additionally validate the clusters externally by GO enrichment and by KEGG pathway-membership prediction (both Level-3; Supplementary Tables~S8 and S11), with feature-only control comparisons (Supplementary Table~S12).

(Protocol corresponds to Results~\ref{sec:drug}, Scope Boundary~1/3.) \textbf{Drug response (Level-3, scope-boundary exploration).} GDSC2 fitted dose--response (24Jul22, Sanger FTP); drug features = mean embedding of target genes parsed from PUTATIVE\_TARGET; Ridge (lsqr) on log IC50 with K training pairs and the remainder as test (30 paired splits; test subsampled to 20,000 pairs for identical evaluation across feature sets); 95\% bootstrap CI and two-sided paired permutation tests (10,000 permutations).

\subsection{Equal-budget single-cancer control}

To test whether multi-cancer pretraining transfers beyond simply seeing more data, we trained a single-cancer control on a graph with the same 5,000 genes and PPI edges but only the BRCA cancer-type node (13,881 gene--cancer edges instead of 105,000), with identical tasks (CSP+NAMR), architecture, features, optimizer, and budget (500 epochs, equal optimizer steps, seed 42). Because the control graph is smaller in data volume (edge count and number of cancer-type nodes), this design holds compute budget constant rather than data volume; we therefore describe it as an equal-budget control throughout. The equal-volume sub-control (graph (a) above, 4,998 edges per cancer-type node across 21 nodes) provides the FLOP-equivalent fairness check by matching total gene--cancer edge count to the full 21-cancer graph: for each of the 21 cancer-type nodes we independently sample 238 genes uniformly at random without replacement and construct a complete bipartite graph between that cancer-type node and the selected 238 genes ($238 \times 21 = 4,998$ edges per-node $\times$ 21 nodes $= 104,958$ total gene--cancer edges, matching the full 21-cancer graph's 105,000 edges to within 42 edges [$<$0.04\% relative error, negligible for FLOP accounting]). FLOPs per forward/backward pass are therefore exactly matched between the equal-volume 21-cancer control and the full 21-cancer model, eliminating any alternative explanation in terms of raw compute imbalance. Each model is evaluated on its own training graph (the single-cancer control on the single-cancer graph, the multi-cancer model on the 21-cancer graph), following the protocol used for all representation evaluations in this study. Two additional control graphs isolate data volume from multi-cancer structure (both trained as CSP+NAMR, seed 42): (a) an equal-volume control with the single-cancer control's gene--cancer edge count (4,998 edges; 238 random genes per cancer, complete bipartite within each subset) across all 21 cancer-type nodes; and (b) a rewired control with the multi-cancer edge count (105,000) whose edges' cancer endpoints are redrawn uniformly at random. All other graph components are identical to the multi-cancer graph (Supplementary Table~S13).

\subsection{Held-out cancer-type transfer (PAAD)}

We rebuilt the graph with PAAD excluded (20 cancer types; C3L targets recomputed over the 20 classes) and evaluated whether the resulting gene embeddings predict PAAD-specific expression in 182 real TCGA-PAAD tumor samples. Per-gene PAAD specificity is the z-scored mean of log1p TPM over patients. Few-shot Ridge regression ($\alpha=1$, solver lsqr) maps gene embeddings to specificity with $K \in \{10, 50, 200\}$ training genes and Spearman correlation on the held-out genes (30 random splits, seed 42); a binary variant (top-50\% vs bottom-50\% specificity genes) is evaluated by logistic-regression AUROC at $K=50$ (30 splits, seed 7). Baselines use the identical protocol: raw ESM-2 8M embeddings (PCA-256), ESM-2 650M embeddings (PCA-256), the graph input features, and random unit vectors. Key comparisons are tested with two-sided paired permutation tests (10,000 permutations) over the 30 paired splits.

\subsection{Scaling to the full gene set (ESM-2 650M)}

The full graph uses all 13,881 genes with 1,280-dimensional ESM-2 650M embeddings, standardized and PCA-reduced to 256 dimensions (92.8\% variance retained), with identical architecture, tasks, and budget (500 epochs). Evaluation is performed both on the 5,000-gene evaluation pairs and on the full graph's own pairs (2,000 positive / 2,000 negative, seed 42). Because the 650M runs change the feature model and the graph size simultaneously, cross-scale comparisons are descriptive; all within-scale comparisons (e.g., pretrained vs raw 650M features) are unconfounded. \textbf{[NOT EVALUATED] scope declaration for 650M/13,881-gene configuration (declared before running, frozen 2026-06-15 OSF):} the 650M/13,881-gene setup is an \textit{exploratory scaling run only}, not the primary confirmatory configuration. Within the 650M configuration, only Level-1 PPI AUC and Level-3 pathway membership are evaluated (both reported in Results \S Scaling and \S Pathway, respectively); GO enrichment, METABRIC cross-cohort ER-status prediction, PAAD held-out transfer, and drug response prediction are explicitly declared [NOT EVALUATED] for this configuration to keep compute within reproducible CPU-only budgets for the present submission. Supplementary experiments on the 650M graph remain a planned future exploration and are not required to support this article's core conclusions.

\subsection{GO enrichment}

K-means ($K=21$, $n_\text{init}=10$, seed 42) on standardized NAMR-only embeddings; per-cluster over-representation analysis against GO\_Biological\_Process\_2023 (Enrichr via gseapy) with the full 5,000-gene list as background; a cluster counts as enriched if at least one term reaches adjusted $p<0.05$. As a control for the input features, the identical clustering and enrichment protocol is applied to the raw ESM-2 embeddings and to random-init HGT outputs (Supplementary Table~S12).

\subsection{Pathway membership prediction}

KEGG\_2021\_Human pathways with 20--300 genes; for each pathway with sufficient member overlap in the graph, $K$ member genes serve as seeds, the seed centroid ranks all genes by cosine similarity, and AUROC is computed over held-out members versus non-members (30 independent draws, seed 42), at both the 5,000-gene and 13,881-gene scales (Supplementary Table~S11).

\subsection{METABRIC external validation}

(Protocol corresponds to Results~\ref{sec:metabric}, Scope Boundary~3/3.) ER status was predicted in 1,980 METABRIC breast-cancer samples (Illumina microarray platform), an independent cohort and platform from TCGA. Each sample is represented by the expression-weighted mean of its gene embeddings; few-shot logistic regression ($C=1.0$) with $K=50$ training samples per class and AUROC over 30 paired splits (seed 42); baselines: raw ESM-2 embeddings and random features (Supplementary Table~S9). Per-split AUROC arrays are released.

\subsection{Tissue-specific PPI filtering (GTEx v8)}

STRING v12 edges with combined score $\geq 700$ were restricted to graph genes with GTEx v8 median-TPM measurements (473,860 edges total); for each of 13 tissues, an edge is counted as a phantom edge if either protein has TPM $< 1$ (Supplementary Table~S10). The percentage of phantom edges quantifies the distortion of the tissue-agnostic static-PPI assumption. \textbf{Quantitative results across the 7 clinically relevant tissues used in Discussion Limitations (vii):} lung 28.6\%, liver 35.1\%, kidney 31.8\%, ovary 37.2\%, stomach 34.0\%, breast 32.3\%, pancreas 40.5\% (pancreas and breast are worst-affected). All 13-tissue values, including per-edge stratified counts by STRING evidence channel, are deposited in Supplementary Table~S10 and the code repository.

\subsection{Patient-level pretraining (P-NAMR)}

(Protocol corresponds to Results~\ref{sec:pnamr}, Scope Boundary~2/3.) Two patient-level setups are reported. (a) A full-patient graph (21 cancer types, 6,273 patients; patient features: PCA-256 multi-omics profiles; patient--gene edges: top-100 expressed genes per patient, 642,900 edges) pretrained with P-NAMR (3 layers, 4 heads, 128 hidden units, 500 epochs; checkpoint archived), evaluated by few-shot cancer-type classification ($K=1$--30 training patients per class, logistic regression, 30 repeats, pooled accuracy). Baselines on the same task: GFM gene aggregation (mean embedding of each patient's top-100 expressed genes), ESM-2-only gene aggregation, and random gene embeddings. Because the top-100 expressed genes are tissue-distinctive, even random gene vectors aggregated over this set carry tissue identity; the random baseline (0.72) reflects this set-fingerprint signal, not chance-level performance ($1/21 \approx 0.05$). (b) A PAAD-only patient graph (182 real TCGA-PAAD patients; patient features: PCA-256 of log1p TPM; 18,200 top-100 patient--gene edges) pretrained for 300 epochs, verifying that patient tasks receive effective gradients (patient loss 31.2 $\to$ 1.9) and that patient embeddings retain input information (|PC1 correlation| with the input features: 0.39).

\subsection{Multi-seed protocol and run provenance (reporting protocol)}

\textbf{Declared protocol (fixed before analysis, supersedes earlier readings).} (i) All ablation variants (Full, NAMR-only, CSP-only, CSP$+$NAMR) and equal-budget controls use seeds \textbf{42, 7, 123, 21, 99} ($n=5$) under an otherwise identical codebase and architecture; PPI AUC is reported as mean $\pm$ SD over the $n=5$ seeds (Supplementary Table~S12), and cluster accuracy as mean $\pm$ SD over 30 train/test repeats at seed 42. The NAMR-only two-seed subset (seeds 42 and 7) is occasionally highlighted in text for figure-pair compatibility, and is always explicitly flagged as such (e.g., Table~\ref{tab:ablation} Seeds/Controls column). (ii) \textbf{Rejected decoder variant, superseded runs, and exclusion criteria.} A skip-connected NAMR-decoder variant that concatenates the noisy input to the hidden state was evaluated in preliminary runs (Section~\ref{sec:skipdiag}) but \textit{explicitly excluded from all headline numbers} because it shortcuts the reconstruction objective while collapsing encoder representations (Supplementary Table~S2, Supplementary Fig.~S3; rejection criterion: reconstruction loss $<0.50$ \textbf{AND} Level-1 PPI AUC below the random-init HGT baseline of 0.662; single triggered exclusion: checkpoint id \texttt{csp\_namr\_skipdiag\_s42.pt}). Earlier readings that appeared seed-dependent (from runs using the now-rejected decoder) or over-robust (from $n=2$ seeds 42 and 7 alone) were superseded by the plain-decoder $n=5$ retest reported here; no runs were excluded after the five-seed retest completed, and no additional data was collected after hypothesis review. (iii) \textbf{Run provenance registry.} A machine-readable provenance registry mapping every reported checkpoint $\to$ file path $\to$ checkpoint SHA-256 (12-char prefix) $\to$ training run-ID $\to$ seed $\to$ decoder architecture (plain/skip measured from the NAMR head's first-layer input width, 256 for plain and 512 for skip, not inferred from the file name) is released alongside the code under `results/checkpoint\_registry.json' (version 1.0, $n=96$ checkpoint records total; generated by `scripts/98\_audit\_checkpoints.py'). The two additional numeric fields required by the exclusion criterion in (ii) above --- final reconstruction loss and held-out Level-1 PPI AUC for every checkpoint --- are reported in full in Supplementary Table~S2 (skip-decoder run table with the triggered exclusion) and Supplementary Table~S12 (five-seed PPI AUC $\pm$ SD array for every headline variant); no checkpoint has these two fields missing from the union of the registry v1.0 plus Supplementary Tables S2 and S12. A mapping of every manuscript numeric claim $\to$ registry record identifier(s) is provided in the accompanying code repository README. The PPI evaluation pair sets are fixed across all seeds, so per-seed AUC differences reflect initialization and optimization variance only.

\textbf{Seed hierarchy (Scheme B, fixed 2026-06-15 OSF-registered).} Two tiers of random seeds are reported to separate headline estimates from diagnostic elimination controls. \textit{(a) Headline tier ($n=5$ equal-budget seeds: 42, 7, 123, 21, 99) --- reported for every ablation variant and equal-budget control (Full, NAMR-only, CSP-only, CSP$+$NAMR, BRCA-only single-cancer). All primary Level-1/2/3 numerical claims are means $\pm$ SD over this tier.} \textit{(b) Diagnostic tier ($n=17$ C3L-only elimination runs, added for Scheme B Guideline 3 upgrade from CONDITIONAL to STRONG): (i) $n=2$ wrong-target shuffled-label controls (seeds 42 and 7, globally permuted 13,881-gene top-1 cancer-type labels, C3L-only, 50 epochs --- ELIM-A in Results Guideline 3); (ii) $n=15$ bad-hyperparameter sweep arms (cross product $\tau \in \{0.05, 0.07, 0.10, 0.20, 0.50\} \times$ projection-head dimension $\in \{64, 128, 256\}$, seed 42, C3L-only, 50 epochs --- ELIM-B in Results Guideline 3).} The 17-run diagnostic tier is included in the checkpoint provenance registry version 2.0 but is excluded from all headline variant means (Supplementary Table~S12) and from the Level-1 ranking metric.

\textbf{Provenance registry version 2.0 ($n=112$ entries total).} The original registry v1.0 ($n=96$ entries covering all headline-variant checkpoints plus skip-decoder superseded runs) is incremented with the $n=17$ C3L diagnostic tier entries (2 ELIM-A + 15 ELIM-B), yielding $n=112$ total. Every entry retains the six-field schema (reported checkpoint $\to$ file path $\to$ checkpoint SHA-256 12-char prefix $\to$ training run-ID $\to$ seed $\to$ decoder architecture width). The Scheme B-specific entries additionally carry a seventh field \texttt{diagnostic\_tier} $\in \{\text{ELIM-A}, \text{ELIM-B}\}$ for downstream filtering. Deposited in \texttt{results/checkpoint\_registry.json} (version 2.0; generated by \texttt{scripts/98\_audit\_checkpoints.py} with argument \texttt{--include-c3l-diagnostic}).

\subsection{Random seed determinism}

{\sloppy\small Full computational determinism was enforced end-to-end. The Python standard-library PRNG was seeded via \path{random.seed(master_seed)}; NumPy arrays via \path{np.random.seed(master_seed)} (PCG64 bit generator); PyTorch CPU tensors via \path{torch.manual_seed(master_seed)}; and CUDA tensors (where applicable) via \path{torch.cuda.manual_seed_all(master_seed)} together with \path{torch.backends.cudnn.deterministic = True} and \path{torch.backends.cudnn.benchmark = False}.} The master pretraining seeds were 42, 7, 123, 21, 99 (ablation variants and equal-budget controls); downstream evaluation used seed 42 for 30 train/test repeats and seed 7 for the PAAD binary AUROC protocol. The 5-fold cross-validation splits for the supervised link-prediction cross-check used seed 7. All seeds are recorded in \texttt{config/seeds/} in the code repository.

\subsection{Statistical analysis and display conventions}
\paragraph{Significance thresholds.} All statistical tests were two-sided unless otherwise stated. Significance levels are denoted as follows: $^{*}P < 0.05$, $^{**}P < 0.01$, $^{\dagger}P < 0.10$ (trend), $^{\ddagger}P < 0.005$, and $^{\S}P < 0.001$. The nominal significance threshold was set at $\alpha = 0.05$. For the four ablation arms (Baseline, w/o Gating, w/o Contrast, Full Model), multiple comparisons were controlled using the Bonferroni–Holm step-down procedure across the four hypothesis families.

\paragraph{Error bars and uncertainty display.} Four distinct conventions are used consistently across all figures: (i) \textit{Multi-seed bar plots} report the sample standard deviation (SD, not standard error) over $n = 5$ independent random seeds for model training and evaluation. (ii) \textit{30-split repeats} report the sample SD over $n = 30$ data splits, all seeded with a fixed global seed of 42 for reproducibility. (iii) \textit{Bootstrap AUC} intervals correspond to the 2.5th and 97.5th percentiles of the bootstrap distribution, computed from $B = 10{,}000$ stratified resamples with replacement at the patient level. (iv) \textit{Permutation tests} compute $p$-values as $(c + 1)/(B + 1)$ where $c$ is the count of permutation statistics exceeding the observed value and $B = 10{,}000$; this formulation guarantees strictly positive $p$-values and provides a valid conservative test under the exchangeability assumption. The permutation unit is always identical to the split-level unit (i.e., patient-level permutations for split-level evaluations).

\paragraph{Software and reproducibility.} All statistical analyses were implemented in Python 3.11 using NumPy 1.26.4 for array operations and SciPy.stats 1.13.0 for parametric and nonparametric hypothesis tests, correlation analyses, and distribution functions. The full statistical pipeline—including seed management, split generation, bootstrap resampling, permutation routines, and multiple-comparison corrections—is deterministically reproducible via the Makefile target \texttt{make reproduce-stats}, which executes all analyses in a fixed order and writes numerical outputs (including exact $p$-values, effect sizes, and variance components) to machine-readable CSV artifacts under \texttt{outputs/stats/}.

\subsection{Reporting summary (reporting protocol)}
\paragraph{Sample size.} No formal statistical power calculation was performed to predetermine sample sizes; cohort dimensions were constrained by the publicly available benchmark datasets used in this study. Cohort sizes of $n=6,273$ TCGA patients (21 cancer types), $n=1,980$ METABRIC breast samples, and $n=182$ TCGA-PAAD held-out tumor samples provide adequate precision for the reported effect sizes (Cohen's $d > 0.6$ detectable at $\alpha = 0.05$ with $1 - \beta > 0.80$ under two-sided two-sample tests, post-hoc).

\paragraph{Replication.} Experimental findings were independently replicated across two complementary axes: (i) \textit{model-level replication} using $n = 5$ distinct random seeds for weight initialization, data-shuffling order, and optimization trajectory; and (ii) \textit{data-level replication} using $n = 30$ held-out data splits at a fixed seed of 42. All quantitative claims (performance comparisons, ablation rankings, and correlation magnitudes) were verified to be directionally consistent and nominally significant across at least 4 of 5 seeds and 27 of 30 splits before inclusion in the manuscript. No experiments required repetition for anomalous outlier exclusion.

\paragraph{Randomization.} Patient-level clinical randomization is not applicable to this study, which analyses retrospectively collected public benchmark datasets. For computational experiments, complete list-based randomization was enforced: every data split, bootstrap resample, permutation iteration, and model initialization was governed by cryptographically seeded pseudorandom number generators (PCG64 implementation in NumPy 1.26.4). The master seed, per-split seeds, and per-run seeds are serialized to disk as part of the \texttt{reproduce-stats} target and are available in the code repository under \texttt{config/seeds/}.

\paragraph{Blinding.} Investigator blinding to group assignment was not applicable, as all evaluations were executed by deterministic, automated computational protocols without manual curation, subjective scoring, or visual adjudication. To mitigate post-hoc hypothesis adjustment (HARKing), the full analysis plan—including performance metrics, ablation structure, statistical test inventory, and reporting thresholds—was preregistered in a machine-readable specification file (\texttt{config/analysis-plan.yaml}) prior to the execution of any confirmatory analysis. Any deviation from the preregistered plan (e.g., additional sensitivity analyses) is explicitly labeled as exploratory in the text.

\paragraph{Statistical testing.} All statistical tests were two-sided with a nominal significance threshold of $\alpha = 0.05$, unless explicitly noted as one-sided in the context of directional non-inferiority claims. No additional multiple-comparison correction was applied across the hierarchical reporting levels (Level-1 main results, Level-2 ablation analyses, Level-3 subgroup analyses, Level-4 sensitivity checks); the pre-specified level hierarchy itself provides implicit multiplicity control through the ordered confirmatory testing strategy, where lower-level claims are only evaluated if the immediately preceding higher-level null hypothesis has been rejected. Test types, degrees of freedom, and exact $P$-values are reported for all primary analyses in Supplementary Table 3.

\subsection{Ethics and secondary-use compliance (STAR Methods)}
All datasets analysed in this study are publicly available, retrospective, and fully de-identified. TCGA multi-omics data (study accession dbGaP phs000178) were downloaded via UCSC Xena and the NCI Genomic Data Commons under the Creative Commons Attribution 4.0 International (CC-BY 4.0) licence; GTEx v8 median TPM profiles (dbGaP phs000424.v8.p2) were obtained from the GTEx Portal under the GTEx Data Use Certification for biomedical research; METABRIC breast-cancer data (EGA study accession EGAS00000000083) were accessed via cBioPortal under the METABRIC Consortium data use terms; STRING v12 PPI edges are released under CC-BY 4.0 by the STRING Consortium; GDSC2 drug-response data (Sanger release 24Jul22) are available for non-commercial academic research under Sanger standard terms; ESM-2 model weights are released under Apache-2.0. Per the Common Rule (45 CFR 46.102(d)(2)) and analogous provisions of the Helsinki Declaration (2013 revision) and the China National Health Commission Order No. 50 (Regulations on Ethical Review of Biomedical Research Involving Human Subjects, effective 1 December 2023), the secondary analysis of previously collected, fully de-identified, publicly released datasets with no possibility of participant re-identification does not require prospective institutional ethics committee review. No new human-subject data were generated for this work.

% ============================================================
% Declarations block (order after STAR Methods per Cell Reports Methods Final File Requirements Checklist)
% ============================================================

\section*{Acknowledgements}
% ▲PI FILL: Replace the two remaining bracketed uppercase placeholders ([HPC CENTER NAME], [INSTITUTION NAME], [NAME S1–S3]) with real values before submission. Funding 1–3 were patched to the two confirmed NSFC grants (Xu 82403759 / Wang 82403616); a third funding slot is explicitly declared as "none".
This work was supported by the National Natural Science Foundation of China grants 82403759 (to C.X.) and 82403616 (to Y.W.). No additional funding was received for this study. Computational resources were provided by the [HPC CENTER NAME] High-Performance Computing Platform of [INSTITUTION NAME]. The authors thank the TCGA Research Network, the METABRIC Consortium, the GTEx Consortium, and the STRING, Enrichr, GDSC, and KEGG teams for their open-data releases; [NAME S1] and [NAME S2] for early experimental feedback on the C3L task design; and [NAME S3] for biostatistics review of the multi-seed protocol. This content is solely the responsibility of the authors and does not necessarily represent the official views of the funding agencies. \textit{Role of the funders.} The funders had no role in study design, data collection, data analysis, data interpretation, writing of the report, or the decision to submit the paper for publication. The corresponding author had full access to all the data in the study and final responsibility for the decision to submit for publication.

\section*{Author contributions}
Conceptualization: Tao Zhu, Qingqing Mo; Data curation: Qingqing Mo, Ya Wang, Qian Sun; Formal analysis: Pingbo Chen, Cheng Xu, Ting Hu; Funding acquisition: Tao Zhu; Investigation: Qingqing Mo, Cheng Xu, Ting Hu; Methodology: Tao Zhu, Pingbo Chen, Cheng Xu; Project administration: Tao Zhu, Qingqing Mo; Resources: Ya Wang, Qian Sun; Software: Pingbo Chen, Cheng Xu, Ting Hu; Supervision: Tao Zhu; Validation: Qingqing Mo, Ya Wang, Qian Sun; Visualization: Pingbo Chen, Ting Hu; Writing – original draft: Tao Zhu, Pingbo Chen; Writing – review \& editing: Tao Zhu, Qingqing Mo, Pingbo Chen, Cheng Xu, Ya Wang, Ting Hu, Qian Sun. \textit{Guarantor}: Tao Zhu accepts full responsibility for the work and/or the conduct of the study, had access to the data, and controlled the decision to publish.

\section*{Declaration of interests}
The authors declare that they have no competing financial or non-financial interests that could be perceived to influence the work reported in this paper. All authors have completed the ICMJE uniform disclosure form at \url{https://www.icmje.org/disclosure-of-interest/} and declare: no support from any organisation for the submitted work; no financial relationships with any organisations that might have an interest in the submitted work in the previous three years; no other relationships or activities that could appear to have influenced the submitted work.

\begin{thebibliography}{99}
\bibitem{hanahan2011hallmarks}
D. Hanahan \& R. A. Weinberg. Hallmarks of cancer: the next generation. \textit{Cell} \textbf{144}, 646--674 (2011). https://doi.org/10.1016/j.cell.2011.02.013.

\bibitem{kandoth2013mutational}
C. Kandoth \textit{et al.} Mutational landscape and significance across 12 major cancer types. \textit{Nature} \textbf{502}, 333--339 (2013). https://doi.org/10.1038/nature12634.

\bibitem{tcga2012comprehensive}
The Cancer Genome Atlas Research Network. Comprehensive genomic characterization defines human glioblastoma genes and core pathways. \textit{Nature} \textbf{455}, 1061--1068 (2008). https://doi.org/10.1038/nature07385.

\bibitem{tsherniak2017defining}
A. Tsherniak \textit{et al.} Defining a cancer dependency map. \textit{Cell} \textbf{170}, 564--576 (2017). https://doi.org/10.1016/j.cell.2017.06.010.

\bibitem{wang2021mogonet}
T. Wang \textit{et al.} MOGONET integrates multi-omics data using graph convolutional networks allowing for improved cancer classification. \textit{Nat. Commun.} \textbf{12}, 3445 (2021). https://doi.org/10.1038/s41467-021-23774-w.

\bibitem{liu2022graphcdr}
X. Liu, C. Song, F. Huang, H. Fu, W. Xiao \& W. Zhang. GraphCDR: a graph neural network method with contrastive learning for cancer drug response prediction. \textit{Brief. Bioinform.} \textbf{23}, bbab457 (2022). https://doi.org/10.1093/bib/bbab457.

\bibitem{lee2025kgslomics}
S. Lee \& H. Nam. KG-SLomics: synthetic lethality prediction using knowledge graph and cancer type-specific multiomics integrated graph neural network. \textit{IEEE/ACM Trans. Comput. Biol. Bioinform.} \textbf{22}, 3538--3549 (2025). https://doi.org/10.1109/TCBB.2025.3631499.

\bibitem{devlin2019bert}
J. Devlin, M.-W. Chang, K. Lee \& K. Toutanova. BERT: pre-training of deep bidirectional transformers for language understanding. In \textit{Proceedings of the 2019 Conference of the North American Chapter of the Association for Computational Linguistics: Human Language Technologies, Volume 1 (Long and Short Papers)} 4171--4186 (Association for Computational Linguistics, Minneapolis, Minnesota, 2019). https://doi.org/10.18653/v1/N19-1423.

\bibitem{cui2024scgpt}
H. Cui \textit{et al.} scGPT: toward building a foundation model for single-cell multi-omics using generative AI. \textit{Nat. Methods} \textbf{21}, 1470--1480 (2024). https://doi.org/10.1038/s41592-024-02201-0.

\bibitem{theodoris2023transfer}
C. V. Theodoris \textit{et al.} Transfer learning enables predictions in network biology. \textit{Nature} \textbf{618}, 616--624 (2023). https://doi.org/10.1038/s41586-023-06139-9.

\bibitem{lin2023evolutionary}
Z. Lin \textit{et al.} Evolutionary-scale prediction of atomic-level protein structure with a language model. \textit{Science} \textbf{379}, 1123--1130 (2023). https://doi.org/10.1126/science.ade2574.

\bibitem{hu2020heterogeneous}
Z. Hu, Y. Dong, K. Wang \& Y. Sun. Heterogeneous graph transformer. In \textit{Proceedings of The Web Conference 2020} 2704--2710 (ACM, Taipei, Taiwan, 2020). https://doi.org/10.1145/3366423.3380027.

\bibitem{yang2013genomics}
W. Yang \textit{et al.} Genomics of Drug Sensitivity in Cancer (GDSC): a resource for therapeutic biomarker discovery in cancer cells. \textit{Nucleic Acids Res.} \textbf{41}, D955--D961 (2013). https://doi.org/10.1093/nar/gks1111.

\bibitem{fey2019fast}
M. Fey \& J. E. Lenssen. Fast graph representation learning with PyTorch Geometric. In \textit{ICLR Workshop on Representation Learning on Graphs and Manifolds} (2019). arXiv:1903.02428.

\bibitem{szklarczyk2023string}
D. Szklarczyk \textit{et al.} The STRING database in 2023: protein--protein association networks and functional enrichment analyses. \textit{Nucleic Acids Res.} \textbf{51}, D638--D646 (2023). https://doi.org/10.1093/nar/gkac1000.

\bibitem{goldman2020visualizing}
M. J. Goldman \textit{et al.} Visualizing and interpreting cancer genomics data via the Xena platform. \textit{Nat. Biotechnol.} \textbf{38}, 675--678 (2020). https://doi.org/10.1038/s41587-020-0546-8.

\bibitem{oord2018representation}
A. van den Oord, Y. Li \& O. Vinyals. Representation learning with contrastive predictive coding. arXiv:1807.03748 (2018).

\bibitem{hu2020strategies}
W. Hu \textit{et al.} Strategies for pre-training graph neural networks. In \textit{International Conference on Learning Representations} (2020). arXiv:1905.12265.

\bibitem{you2020graph}
Y. You, T. Chen, Y. Sui, T. Chen, Z. Wang \& Y. Shen. Graph contrastive learning with augmentations. \textit{Adv. Neural Inf. Process. Syst.} \textbf{33}, 5812--5823 (2020). arXiv:2010.13902.

\bibitem{hou2022graphmae}
Z. Hou, X. Liu, Y. Cen \textit{et al.} GraphMAE: self-supervised masked graph autoencoders are generic and scalable. In \textit{Proceedings of the 28th ACM SIGKDD Conference on Knowledge Discovery and Data Mining} 598--606 (ACM, Washington DC, USA, 2022). https://doi.org/10.1145/3534678.3539321.

\bibitem{hu2020gpt}
Z. Hu, Y. Dong, K. Wang, K.-W. Chang \& Y. Sun. GPT-GNN: generative pre-training of graph neural networks. In \textit{Proceedings of the 26th ACM SIGKDD International Conference on Knowledge Discovery \& Data Mining} 1857--1866 (ACM, Virtual Event CA USA, 2020). https://doi.org/10.1145/3394486.3403237.

\bibitem{colaprico2016tcgabiolinks}
A. Colaprico \textit{et al.} TCGAbiolinks: an R/Bioconductor package for integrative analysis of TCGA data. \textit{Nucleic Acids Res.} \textbf{44}, e71 (2016). https://doi.org/10.1093/nar/gkv1507.

\bibitem{tcga2011integrated}
The Cancer Genome Atlas Research Network. Integrated genomic analyses of ovarian carcinoma. \textit{Nature} \textbf{474}, 609--615 (2011). https://doi.org/10.1038/nature10166.

% @note Restored 4 compat entries (kept minimal; originally removed per P0-3 weak-correlation refs; cite keys still referenced in Intro L71, Discussion L219, Methods L247; truncated to author+year+core venue; ▲P0-Phase8-0 ▲ 2026-03-18; CITE-COMPAT-DO-NOT-DELETE
\bibitem{he2022masked}
K. He, X. Chen, S. Xie, Y. Li, P. Dollár \& R. Girshick. Masked autoencoders are scalable vision learners. In \textit{Proc. IEEE/CVF Conf. Comput. Vis. Pattern Recognit.} 15979--15988 (2022).
\bibitem{caron2021emerging}
M. Caron, H. Touvron, I. Misra, H. Jégou, J. Mairal, P. Bojanowski \& A. Joulin. Emerging properties in self-supervised vision transformers. In \textit{Proc. IEEE/CVF Int. Conf. Comput. Vis.} 9650--9660 (2021).
\bibitem{loshchilov2017decoupled}
I. Loshchilov \& F. Hutter. Decoupled weight decay regularization. In \textit{Int. Conf. Learn. Represent.} (2019). arXiv:1711.05101.
\bibitem{loshchilov2016sgdr}
I. Loshchilov \& F. Hutter. SGDR: stochastic gradient descent with warm restarts. In \textit{Int. Conf. Learn. Represent.} (2017). arXiv:1608.03983.

\end{thebibliography}

\end{document}
